\documentclass[11pt,letterpaper,reqno]{amsart}
% BEGIN INLINE manuscript-style
\makeatletter
\NeedsTeXFormat{LaTeX2e}[2020/10/01]
% Use with \documentclass[11pt,letterpaper,reqno]{amsart}.
% This file controls presentation only. It supplies no mathematical content.

\RequirePackage[T1]{fontenc}
\RequirePackage{lmodern}
\RequirePackage{microtype}
\RequirePackage[letterpaper,margin=1in,headheight=15pt]{geometry}
\RequirePackage{amsmath,amsthm,amssymb,mathtools}
\RequirePackage{enumitem}
\RequirePackage{perpage}
\RequirePackage{longtable}
\RequirePackage{xcolor}
\definecolor{manuscriptred}{rgb}{0.6,0.1,0.1}
\definecolor{manuscriptblue}{rgb}{0.1,0.3,0.5}
\definecolor{manuscriptmagenta}{HTML}{60445C}
\RequirePackage[breaklinks=true,colorlinks=true,
  linkcolor=manuscriptred,citecolor=manuscriptblue,urlcolor=manuscriptmagenta]{hyperref}
\RequirePackage{cleveref}
\hypersetup{pdfauthor={}}

% Retain the AMS paragraph, title, heading and running-head conventions.
% Equation counters restart each section: 2.i, 2.ii, ... .
% Appendix sections keep their letters: A.i, A.ii, ... .
\numberwithin{equation}{section}
\renewcommand{\theequation}{\thesection.\roman{equation}}
% Prose references include "eq." and omit parentheses around the number.
\renewcommand{\eqref}[1]{\textup{eq.~\ref{#1}}}
\crefname{equation}{eq.}{eqs.}
\Crefname{equation}{eq.}{eqs.}
\creflabelformat{equation}{#2#1#3}
\setcounter{tocdepth}{1}
\renewcommand{\thefigure}{\Roman{figure}}
\renewcommand{\thetable}{\Roman{table}}
\newcommand{\emneedspace}[1]{\par\begingroup
  \dimen0=\pagegoal\advance\dimen0 by-\pagetotal
  \ifdim\dimen0<#1\relax\newpage\fi\endgroup}
\MakePerPage{footnote}
\renewcommand{\thefootnote}{%
  \ifnum\value{footnote}>9\arabic{footnote}\else\fnsymbol{footnote}\fi}

% Upright statements. Main results are A, B, C; supporting statements
% share a section-local sequence. These end marks indicate typography,
% never proof completeness, validation, or a deferred-proof status.
\theoremstyle{definition}
\newtheorem{ms@theorem}{Theorem}
\renewcommand{\thems@theorem}{\Alph{ms@theorem}}
\newtheorem*{ms@theoremstar}{Theorem}
\newtheorem*{ms@theoremzero}{Theorem 0}
\newtheorem{ms@proposition}{Proposition}[section]
\newtheorem{ms@lemma}[ms@proposition]{Lemma}
\newtheorem{ms@definition}[ms@proposition]{Definition}
\newtheorem{ms@corollary}[ms@proposition]{Corollary}
\newtheorem*{ms@corollarystar}{Corollary}

\newcommand{\ms@majorstart}{%
  \pushQED{\qed}\renewcommand{\qedsymbol}{$\blacksquare\square$}}
\newcommand{\ms@minorstart}{%
  \pushQED{\qed}\renewcommand{\qedsymbol}{$\blacksquare$}}
\newenvironment{theorem}[1][]{\ms@majorstart\ms@theorem[#1]}
  {\popQED\endms@theorem}
\newenvironment{theorem*}[1][]{\ms@majorstart\ms@theoremstar[#1]}
  {\popQED\endms@theoremstar}
\newenvironment{theoremzero}[1][]{\ms@majorstart\ms@theoremzero[#1]}
  {\popQED\endms@theoremzero}
\newenvironment{proposition}[1][]{\ms@minorstart\ms@proposition[#1]}
  {\popQED\endms@proposition}
\newenvironment{lemma}[1][]{\ms@minorstart\ms@lemma[#1]}
  {\popQED\endms@lemma}
\newenvironment{definition}[1][]{\ms@minorstart\ms@definition[#1]}
  {\popQED\endms@definition}
\newenvironment{corollary}[1][]{\ms@majorstart\ms@corollary[#1]}
  {\popQED\endms@corollary}
\newenvironment{corollary*}[1][]{\ms@majorstart\ms@corollarystar[#1]}
  {\popQED\endms@corollarystar}
% The source's lemma/proposition variants are deliberately not carried over:
% choose one marker for each statement type rather than encode proof status.

\theoremstyle{remark}
\newtheorem{remark}[ms@proposition]{Remark}
\newtheorem*{remark*}{Remark}
% The ordinary amsthm proof environment retains its open square.

\crefname{ms@theorem}{theorem}{theorems}
\Crefname{ms@theorem}{Theorem}{Theorems}
\crefname{ms@proposition}{proposition}{propositions}
\Crefname{ms@proposition}{Proposition}{Propositions}
\crefname{ms@lemma}{lemma}{lemmas}
\Crefname{ms@lemma}{Lemma}{Lemmas}
\crefname{ms@definition}{definition}{definitions}
\Crefname{ms@definition}{Definition}{Definitions}
\crefname{ms@corollary}{corollary}{corollaries}
\Crefname{ms@corollary}{Corollary}{Corollaries}
\crefname{remark}{remark}{remarks}
\Crefname{remark}{Remark}{Remarks}

% Bibliography configuration for this manuscript.
\makeatother
% END INLINE manuscript-style


\title[Iteration of regular interactions]{Iteration of regular interactions on prescribed coupling prefixes}
\date{Working draft, 9 November 2026}
\begin{document}
\raggedbottom
\begin{abstract}
We prove uniform bounds for two chains of regular interactions formed
by prescribed small-field integration steps. The unmarked chain uses
fixed physical radii and no inherited remainders. The full chain
admits older remainders on running radii. A summable age kernel and
a squared-radius remainder sum control their difference uniformly
in the length of the prefix. Physical Cauchy estimates then give the
marked quadratic action coefficient, with its radius loss retained.
For complete profiles with their scalar coefficient subtracted,
an exact transport law gives a uniform forward stability estimate.
The conclusions assume the actual one-step formation, domain and
companion estimates, the stated flow comparisons, and the supplied
inherited remainder bounds. They do not construct those inputs or
solve for a coupling trajectory.
\end{abstract}
\maketitle
\tableofcontents

\section{Introduction}

A local interaction created by block integration is evaluated again
at later scales. Its repeated contribution depends on what has been
subtracted from it. A regular family carries its own quadratic action
subtraction, and the one-step estimate supplies a decaying weight
for its age. A vacuum-subtracted remainder has no such quadratic
subtraction. Its bound instead retains the squared ratio of the
evaluation radius to its creation radius. The distinction matters
when the number of integration steps is unbounded.

We study the iteration on a given finite prefix of reference and
running data. The reference data fix the integration coordinates,
cutoffs and Gaussian conventions. The action weights vary within
their declared complex domains. The two chains use these same
choices. Their common part cancels only after the chain on fixed
radii has been restricted as the same actual function to the
running domain. This identity is stronger than equality of norm
bounds or agreement of separately chosen logarithms.

The proof first makes the required one-step hypotheses explicit.
It then sums the regular age kernel and the inherited remainder
load. A weighted supremum induction bounds the marked difference.
The coefficient readout uses the physical background field, so a
second derivative costs the inverse square of the physical
analyticity radius. Finally, a norm adapted to exact transport
gives forward stability for complete regular profiles. This last
argument controls the profile for a prescribed coupling sequence;
it does not impose the equation determining that sequence.

Conditional small-field effective-action induction is treated in
\cite[Theorems~1 and~3, pp.~259 and~264; hypotheses on pp.~260--264]{Bal87}.
The local cluster expansion and its effective-action estimates are
completed in \cite[\S2, pp.~11--22]{Bal88c}.
The arguments below give the finite-prefix summations, domain
compatibility and coefficient estimates in the explicitly stated
two-chain setting. The marked recurrence is not attributed to
those cited theorems.

The inherited remainder estimates, their analytic domains, and the
actual one-step theorem are premises of this manuscript. The
quantitative statements distinguish the implications proved from
these premises from their admission in a full lattice construction.
No boundary or large-field history estimate is inferred from a
regular-family norm.

\section{The regular data at a prescribed step}
\label{em:sec-regular-data}

Fix an odd block factor $L\ge3$ and a localization side $M$, a power of
$L$. Use centered compatible grids on the finite four-dimensional tori
and their unit-lattice covers. A level-$j$ domain $X$ is a connected union
of level-$j$ localization cubes. In level-$j$ units let $d_j(X)$ be the
length of a shortest tree in $X$ meeting every cube, divided by $M$.
We use the actual determining windows of the step, including the collars
needed for the fluctuation response and for the pointing at unit sites
$z\in X$. Every assertion about a common lift includes these windows.
Unless a result states a more specific dependence, an unsubscripted
$C$ may change from line to line and depends only on the fixed
geometry, rates and admitted domain bounds. It is independent of
the prefix length, fine mesh and compatible volume. Named input
constants remain fixed when the common coupling floor is increased.

The real gauge group is $\operatorname{SU}(2)$ and its complexification
is $\operatorname{SL}_2(\mathbb C)$. Use traceless Hermitian Lie
coordinates with pairing $\langle X,Y\rangle=\tfrac12\operatorname{tr}(XY)$,
continued bilinearly. The Wilson action is
$A(U)=\sum_p(1-\tfrac12\operatorname{tr}U_p)$ over the stored positive
plaquettes, with $U_p$ the oriented plaquette product. The physical
current $J$ is a bond one-form in the units of its level.

Write $P=(\mathbf U,J)$ for a configuration and current pair on a read
window. The current is an independent complex variable in a local core.
A core with radii $\mathbf a=(a_0,a_1)$ uses representatives
$\mathbf U=e^{i\eta_j A'}U$ with $U$ real and $\eta_j$ the bare spacing
in level-$j$ units. The radius $a_0$ bounds the real and complex plaquette
defects in units $\eta_j^2$ and the current in its level-$j$ units.
The radius $a_1$ bounds $A'$ and its covariant first differences.
The real factor must lie in its declared local gauges. The corresponding
curvature and current bounds, with the prescribed changes of units and
margins, are required at every restricted tower used by the step.
Denote the resulting admitted core by
$\mathcal U_j(X;\mathbf a)$. We use the core inclusions supplied by the
background and one-step theorems. In particular, a substitution along a
complex decoupling contour must map into the core of the potential being
read, including its original tower and collar data. A radius comparison
by itself is not a substitution theorem.

A pointed family $G$ at level $j$ assigns a holomorphic function
$G(X,P,z)$ to each $(X,z)$. The value is gauge invariant under the
declared action and depends on the pair, independently of its splitting
into $U,A',J$. These analytic families form complex linear spaces.
Their real subspaces consist of families that take real values on
real admitted pairs. At real values of all input parameters the
supplied input and output families belong to these real subspaces.
Complex directions and complex parameter values are not subject
to that real-valuedness condition.
Its norm is
\begin{equation}\label{em:eq-pointed-norm}
 \|G\|_{j,\kappa;\mathbf a}
   =\sup_{\mathbb T,\,X,\,z\in X}\ 
           e^{\kappa d_j(X)}\sup_{P\in\mathcal U_j(X;\mathbf a)}
                 |G(X,P,z)| .
\end{equation}
Here the supremum includes the admissible tori and the infinite-lattice
companion. An unpointed family obeys the same analyticity, gauge and
reality conventions on its admitted cores and uses the same norm
without $z$. Parameter
suprema are also taken when families are prescribed on a complex
parameter domain. There is no sum over $z$ in this norm. The rooted
domain sums use the separate entropy bound of the one-step construction
at the chosen decay rate. Restricting a family to a smaller core reduces
its norm.

The block map $\mathcal M$ is the prescribed one-step average from
level-$j$ link fields to level-$(j+1)$ link fields, expressed in their
respective units. Its iterates are compositions of these actual maps.
For a level-$j$ datum $W$, let $U_j(W)$ be the supplied fine
constrained background with $\mathcal M^j U_j(W)=W$, and set
$U_0(W)=W$. Its gauge representatives, local analytic continuations
and physical current are part of the admitted background data.
The block and gauge geometry is developed in \cite{BA}; its
nonlinear stationary background is an input here.

A companion $\widehat G$ is one family on the cover that agrees with
each torus function on every common nonwrapping read, after transport
of all determining data. We require this agreement for the same frame
and weights through the forming step and all compatible cutoffs.
Translations preserving the partition act on individual terms.
For the real fields with backgrounds
$P_j(W)=(U_j(W),J(U_j(W)))$ in the required cores, the pointed total is
\[
                   \mathcal P_z(G;W)=\sum_{X\ni z}G(X,P_j(W),z).
\]
The scalar Euclidean class requires
$\mathcal P_{rz}(G;rW)=\mathcal P_z(G;W)$ for the admitted lattice
isometries, including unit translations, on the common covariance
class. This is a property for each base point. The invariance of the
sum over base points alone would not give it.

A regular member has this covariance, a companion and its scalar
coefficient $b_j(\widehat G)$. This is the linear readout from the
companion's physical Hessian, defined below on the stated response and
summability domains. Its own Wilson counterterm is pointed using the
same midpoint tents as the regular family. We require the actual
fine Wilson action $A(U_j(W))$ to have coefficient one in that
readout. Unit-translation covariance of the curl response and
affine-curl reproduction by the supplied background give this normalization by the
calculation in \cref{em:sec-coefficient}; their admission is part of
the background input. In level-$j$ units, put
\begin{equation}\label{em:eq-tents}
                h_z(x)=\prod_{\mu=1}^4(1-|x_\mu-z_\mu|)_+.
\end{equation}
Periodize these functions on a torus. The two neighboring linear
weights in each coordinate sum to one, so $\sum_z h_z(x)=1$.
Thus the sum of the counterterms is
$-b_j(\widehat G)A(U)$, with $A$ the Wilson action. Each member is first
vacuum subtracted. For an unpointed remainder $H$ we subtract only its
vacuum value. Its domains are unions of the older cells of side
$M\check R_j$ at its birth scale. It has a companion and covariant terms
under the translations preserving that partition. Its total is invariant
on its declared background class under the allowed partition isometries.


\section{Summing the inherited loads}
\label{em:sec-loads}

Use the fixed blocking factor $L$ and localization side $M$. Fix
an exponent $0<\alpha_{\rm H}<1$. The age of a term of origin
$j$ in the step from $m$ to $m+1$ is $n=m+1-j$. The regular
comparison uses the weight
\begin{equation}\label{em:eq-regular-age}
  \omega_n=(n+1)^{3/2}L^{-\alpha_{\rm H}(n-1)},\qquad n\ge1.
\end{equation}
This is the weight in the one-step estimate after each regular
member's own quadratic action has been subtracted. The use of that
estimate includes its stated read domains and the fixed bound for
subtracting each member's own counterterm. Its numerical load is
charged to the original family norms, as specified below.

\begin{lemma}[Summability of regular ages]\label{em:lem-age-kernel}
The sum $\Omega=\sum_{n\ge1}\omega_n$ is finite. It stays finite
after multiplying the summand by any fixed nonnegative power of
$n+1$. More generally, a supplied kernel bounded by
$C(n+1)^dL^{-\alpha(n-1)}$ is bounded by a fixed multiple of
\eqref{em:eq-regular-age} for any
$0<\alpha_{\rm H}<\min\{1,\alpha\}$. The multiple can depend
on $d,L,\alpha-\alpha_{\rm H}$.
\end{lemma}

\begin{proof}
For each fixed $s\ge0$ the ratio of successive terms of
$(n+1)^s\omega_n$ tends to $L^{-\alpha_{\rm H}}<1$. After a
fixed age these ratios are bounded by a number less than one.
This proves the two summability assertions. For the comparison,
the quotient of the proposed majorant and $\omega_n$ is
$C(n+1)^{d-3/2}L^{-(\alpha-\alpha_{\rm H})(n-1)}$. The positive
exponential reserve bounds this quotient uniformly in $n$.
\end{proof}

For the running-radius chain fix real numbers
$\theta_0,\ldots,\theta_N$ above a floor $X>e$ and constants
$\lambda_\flat>0$, $c_5\ge0$ and $\bar B>0$. Assume the weak
flow comparison
\begin{equation}\label{em:eq-weak-flow}
  \lambda_\flat(j-i)-c_5\le\theta_i-\theta_j
                     \le\bar B(j-i),\qquad 0\le i\le j\le N.
\end{equation}
For example, when $0<\lambda_\flat\le\bar B$, the frame
$\theta_j=X+\lambda_\flat(N-j)$ meets these inequalities for
every finite $N$.
This permits bounded excursions between adjacent entries. The
frame is fixed when the action weights vary on their admitted
domains. It supplies steps $m+1$ for $0\le m\le N$.
For fixed positive $A_1,C_0,C_1,\beta_{\rm s}$ and
$q>\mathfrak p_0\ge0$, use
\begin{align}
  \alpha_\iota(y)&=C_\iota y^{-1/2}(\log y)^q,
          &\delta(y)&=A_1y^{-1/2}(\log y)^{\mathfrak p_0},\nonumber\\
  \bar a_{m,\iota}&=(1+\beta_{\rm s})\alpha_\iota(\theta_m),
          &\check a_{j,\iota}&=\min\{\alpha_\iota(\theta_{j-1}),
                                      \alpha_\iota(\theta_j)\},
  \label{em:eq-running-radii}\\
  \mathsf s_m&=\delta(\theta_m)/\bar a_{m,0},
          &w_{m,j}&=\max_{\iota=0,1}
                       (\bar a_{m,\iota}/\check a_{j,\iota})^2.
  \label{em:eq-running-weight}
\end{align}
The enlarged output radius $\bar a_{m,\iota}$ is fixed before
forming step $m+1$. For an inherited term of origin $j$, the
radius $\check a_{j,\iota}$ is a permitted restriction of its
creation domain. The weight $w_{m,j}$ is the squared ratio
used in the inherited-remainder comparison. Domain admission and
the quadratic evaluation estimate that supplies this ratio belong
to the one-step input. A comparison of numerical radii alone
gives no evaluation identity.

For fixed $r>0$ set
\[
  R(y)=\min\{L^b:b\in\mathbb Z_{\ge0},\ L^b\ge(\log y)^r\},
  \qquad \check R_j=R\bigl(\min_{0\le i\le j}\theta_i\bigr).
\]
An inherited origin $j\le m$ passes the numerical age gate of
step $m+1$ when $M\check R_j\le L^{m+1-j}$. Insertion also
requires the supplied companion and analytic admission for that
term. Write $J_m$ for the designated existing origins passing the
numerical gate, and $J_m^{\rm rec}$ for the existing origins that
fail it. These lists are fixed on each parameter domain, since
the frame and the existing origin labels are fixed there. Every
member of $J_m$ must satisfy its companion and analytic hypotheses
throughout that domain. Failure of such a hypothesis makes the step
inadmissible there; it does not silently remove that member from
the list. Origins in $J_m^{\rm rec}$ are outside the regular exponent
of this step.

\begin{lemma}[Squared-radius loads]\label{em:lem-remainder-load}
Let $\kappa_0>4$ and $p=(\kappa_0-2)/2$. Assume the creation
norms of the inherited remainders are bounded by
$\rho_j\le C_R\theta_j^{-\kappa_0/2}$. For a fixed floor high
enough in terms of the stated constants, one has
\begin{equation}\label{em:eq-aged-load}
  \sum_{j\in J_m} w_{m,j}\rho_j
  \le \sum_{j=1}^m w_{m,j}\rho_j
  \le C_{\mathfrak h}\theta_m^{-p},\qquad
  C_{\mathfrak h}\le C\left(1+\frac{1}{\lambda_\flat(p-1)}\right).
\end{equation}
The constant $C$ may depend on $p,q,\beta_{\rm s},C_R,c_5,\bar B$.
For $1\le j\le m$ the radii also satisfy
\begin{equation}\label{em:eq-radius-nesting}
  \check a_{j,\iota}\le\bar a_{m,\iota}
       \le C(m+2-j)^{1/2}\check a_{j,\iota},\qquad \iota=0,1.
\end{equation}
The numerical gate, applied to the full set of origins, selects
an initial segment. For its recent complement,
\begin{equation}\label{em:eq-recent-load}
  \sum_{j\in J_m^{\rm rec}}w_{m,j}\rho_j
  \le C(1+\log\log\theta_m)^2\theta_m^{-\kappa_0/2}.
\end{equation}
The constant in this last bound can also depend on $r,L,M$.
All bounds are independent of $m,N$ and the frame above the floor.
\end{lemma}

\begin{proof}
The fixed constants $C_\iota$ cancel in each ratio in
\eqref{em:eq-running-weight}. In particular that weight is
\[
 (1+\beta_{\rm s})^2\max_{i\in\{j-1,j\}}
   \frac{\theta_i}{\theta_m}
        \left(\frac{\log\theta_m}{\log\theta_i}\right)^{2q}.
\]
By the lower bound in \eqref{em:eq-weak-flow}, both $\theta_i$
in this maximum are at least $\theta_m-c_5$. Above a fixed floor
the logarithmic ratio is bounded by two. The upper bound for one
step gives $\theta_{j-1}\le\theta_j+\bar B$. Increasing the floor
once more yields
\begin{equation}\label{em:eq-weight-comparison}
    w_{m,j}\le C\theta_j/\theta_m\le C(1+m+1-j).
\end{equation}
In the second inequality use the upper bound from $j$ to $m$.
This also gives the upper inequality in \eqref{em:eq-radius-nesting}.
For the lower one take $X-c_5>e^{2q}$, so each $\alpha_\iota$
decreases on $[X-c_5,\infty)$. From $\theta_j\ge\theta_m-c_5$,
\[
 \check a_{j,\iota}\le \alpha_\iota(\theta_m-c_5)
     \le (1-c_5/\theta_m)^{-1/2}\alpha_\iota(\theta_m).
\]
The last factor in front of $\alpha_\iota$ is at most
$1+\beta_{\rm s}$ after increasing the fixed floor.
The summands in \eqref{em:eq-aged-load} are therefore bounded by
$C\theta_m^{-1}\theta_j^{-p}$. The lower flow comparison gives
$\theta_j\ge (\theta_m+\lambda_\flat(m-j))/2$ if $X\ge2c_5$.
The summand $t\mapsto(\theta_m+\lambda_\flat t)^{-p}$ decreases.
Its sum over nonnegative integers is at most its value at zero
plus its integral on $[0,\infty)$, namely
\[
   \theta_m^{-p}+
                \frac{\theta_m^{1-p}}{\lambda_\flat(p-1)}.
\]
Multiplication by $C\theta_m^{-1}$ proves the bound. The empty
sum at $m=0$ is zero.

For the gate, put $y_j=\min_{i\le j}\theta_i$. If the minimum
drops, $y_j-y_{j+1}\le\theta_j-\theta_{j+1}\le\bar B$.
Above a fixed floor this makes the ratio of
$(\log y_j)^r$ to $(\log y_{j+1})^r$ less than $L$. Rounding up to
the next $L$-power gives $\check R_j/\check R_{j+1}\in\{1,L\}$.
Thus $L^j\check R_j$ is nondecreasing. The inequality defining the
gate is $M L^j\check R_j\le L^{m+1}$, proving the initial-segment
property.

If an origin fails the gate, its age $n=m+1-j$ satisfies
\[
  L^n<M\check R_j\le ML(\log\theta_j)^r
     \le C(\log\theta_m)^r(1+\log(n+1))^r.
\]
The last inequality follows from
$\theta_j\le\theta_m+\bar B(n-1)$ and $\log\theta_m\ge1$.
The last factor is bounded by $C_{L,r}L^{n/2}$ for all $n\ge1$.
Taking logarithms gives $n\le C(1+\log\log\theta_m)$. Each such
origin has $\rho_j\le C\theta_m^{-\kappa_0/2}$ by the lower flow
bound. Sum $C(1+n)$ from \eqref{em:eq-weight-comparison} over these
distinct ages to obtain \eqref{em:eq-recent-load}.
\end{proof}

One common geometric margin also follows numerically from these
profiles. If the supplied curvature threshold is
$\varepsilon(y)=A_0y^{-1/2}(\log y)^{\mathfrak p_0}$ with $A_0>0$,
then
\[
 \frac{\check R_m^2\varepsilon(\theta_m)}{\alpha_0(\theta_m)}
 \le\frac{L^2A_0}{C_0}
                (\log\theta_m)^{\mathfrak p_0+2r-q}.
\]
Indeed the running minimum is at most $\theta_m$, and rounding
gives $\check R_m\le L(\log\theta_m)^r$. Thus the stronger choice
$q>\mathfrak p_0+2r$ makes this ratio uniformly small above a fixed
floor. This supplies that numerical margin when the one-step
geometry requires it. Actual read-domain and collar admission
remain separate hypotheses.

The positive lower slope is essential for the aged-load bound.
For a constant frame the weight is constant and the number of
origins passing the numerical gate grows with the prefix. At
$\kappa_0=4$ even a linearly decreasing frame can give a logarithmic
loss in \eqref{em:eq-aged-load}. For $\theta_{m-n}=X+n$,
$m=X^2$ and $X$ a large integer, take $\rho_j=\theta_j^{-2}$ and
sum over the lags $n=m-j$ from $X$ to $X^2-1$. These origins pass
the gate since $M\check R_j\le C(\log X)^r<L^{n+1}$. The logarithms
of their frame entries are comparable to $\log X$, so their
contribution is bounded below by $cX^{-1}\log X$.


\section{The two chains and their common step}


Let $\theta_0,\ldots,\theta_m$ be the supplied real frame prefix, above a
common floor. Write $x_k$ for the prescribed Wilson weight and
$\tau_k=x_k/\theta_k$. The frame fixes the lattice partitions, Gaussian
references, real fluctuation coordinates and tests. The variable $x_k$
does not change those choices. For the absolute chain fix a pair
$\mathbf a_*$ of small positive radii. Its output and inherited cores
use $\mathbf a_*$ throughout. In the convention that the one-step
theorem enlarges its target radii by $1+\beta_{\rm s}$, its target radii before
enlargement are $\mathbf a_* /(1+\beta_{\rm s})$.

For the running chain use the creation radii $\bar{\mathbf a}_k$
and the inherited radii $\check{\mathbf a}_j$ in
\eqref{em:eq-running-radii}. The latter are
formed only once $\theta_j$ is supplied. The inherited restriction uses
the same function on the smaller core. We assume the admitted core nesting,
including $\bar a_{k,\iota}\le a_{*,\iota}$, on the common parameter
domains. For both systems the radii satisfy the one-step admissibility
conditions, including $CMa_0<a_1$ and $a_1^2\le C'a_0$ with the fixed
chart constants. If a separate current radius is used, it is subject to
the same componentwise domain and ratio tests. The present convention
ties it to $a_0$.

At input step $k$ the full list uses the regular members from levels
$1,\ldots,k$ and the designated aged remainders $H_j$ at the origins
$J_k$ of \eqref{em:eq-aged-load}. The prefix must supply and admit all
its required aged inputs. The inequality $M\check R_j\le L^{k+1-j}$
is their numerical age test.
Since the lengths are powers of $L$, the old cell side then divides
the new unit side in old units. With centered aligned partitions, the
new lattice translations and signed coordinate permutations therefore
preserve the old cell partition. This supplies the geometric part of
the old remainder's symmetry at insertion. Its background membership
still has to hold. A family that fails this age test stays outside the
regular list. The unmarked list uses only its own regular outputs, from
levels $1,\ldots,k$. Both lists begin with no entries.

For either list $\mathfrak f$ let $S_{\mathfrak f}(U)$ be the sum of its
vacuum-subtracted regular families, their own Wilson counterterms and
its vacuum-subtracted remainders. The current used in this expression
on a physical background is its actual current. Each older entry is
evaluated in its own units and on its own determining data. The regular
one-step theorem is an input for the following actual finite integral.
At real admitted data let $K_{k,\theta_k}(W,dV)$ be the fixed all-small-
field kernel on $\mathcal M(V)=W$. It uses the given fibre density,
gauge weight and fluctuation cut. In the resolved fibre chart its real
free array is $B$ and its field is $V'(W,B)$ with $\mathcal M(V')=W$.
The completed fluctuation logarithm includes the pivot links. The
analytic continuations and potential bounds used below refer to that
completed map. The admitted background nesting identifies
$V'(W,0)=\mathcal M^kU_{k+1}(W)$ and
$U_k(V'(W,0))=U_{k+1}(W)$ up to a site gauge. Thus gauge-invariant
carried differences vanish at zero fluctuation. These identities
are required of the supplied backgrounds; their existence is not
deduced from the integral formula. Put
\begin{align}
 \mathcal E_{k+1}[x_k,\mathfrak f](W)
  ={}&\log\!\int K_{k,\theta_k}(W,dV)
          e^{-x_k A(U_k(V))+S_{\mathfrak f}(U_k(V))}\notag\\
     &+x_k A(U_{k+1}(W))-S_{\mathfrak f}(U_{k+1}(W)).
 \label{em:eq-actual-step}
\end{align}
The integral retains its prescribed numerical normalization. The
regular one-step expansion separates a field-independent scalar
$\mathfrak v_{k+1}$ from this physical logarithm:
\[
 \mathcal E_{k+1}[x_k,\mathfrak f](W)
   =\mathfrak v_{k+1}+\sum_{X,z}E_{k+1}(X,P_{k+1}(W),z).
\]
The scalar includes the dimensional and reference-normalizer
constants specified by the actual kernel. At the same frozen step
and action weight this scalar is independent of the incoming list.
The same convention is used on every common read, and this scalar is
retained separately in the exact identity. The regular norm below
bounds the localized remainder in this identity, not an arbitrary
unnormalized logarithm. The scalar has zero curvature coefficient.
For a real physical datum, real action weight and real inherited
data, the logarithm is the real logarithm of the positive integral.

The reference-weight and weight-shifted Gaussian presentations use the
same integral. In its supplied fibre quadratic splitting let
$\Delta_A$ be the action precision, $\Delta_G$ the gauge precision and
$\Delta=\Delta_A+\Delta_G$. The weight shift in the action then obeys
\begin{equation}\label{em:eq-frozen-gaussian-rewrite}
 e^{-\langle B,\Delta B\rangle/2}
 e^{-(\tau_k-1)\langle B,\Delta_A B\rangle/2}
 =e^{-\tau_k\langle B,\Delta B\rangle/2}
  e^{(\tau_k-1)\langle B,\Delta_G B\rangle/2}.
\end{equation}
If $B$ has dimension $d$ and the scaled Gaussian is normalized,
retain the multiplier $\tau_k^{-d/2}$ relative to the normalization
at $\tau_k=1$. Here $\tau_k^{-d/2}=\exp[-(d/2)\log\tau_k]$,
with the logarithm continued from zero at $\tau_k=1$ on the
relative disc $|\tau_k-1|<\sigma_*\le1/2$. The gauge compensation on the
right is still a factor on the cut. The pure Gaussian value of its
integral in the action-only split of \cite{GI} does not replace it in
the interacting integral. On the unrestricted Gaussian with the
admitted action--gauge splitting and residual gauge dimension
$d_{\mathrm g}$, its normalized compensation is
$\tau_k^{d_{\mathrm g}/2}$ by \cite[Proposition~4.3]{GI}. The two
scalar logarithms are therefore $-(d/2)\log\tau_k$ and
$(d_{\mathrm g}/2)\log\tau_k$. In the stated decomposition the
first belongs to $\mathfrak v_{k+1}$, while the second is the
physical total of the Gaussian mismatch family. This unrestricted identity leaves
the compensation factor on the cut unchanged. All remaining
vertices and the cut are kept.

Here are the one-step inputs used by the induction. They concern the
same actual kernel and localization on all stated domains.

\begin{enumerate}
\item The real backgrounds and their complex extensions, the actual
  fibre chart, the completed response and the cuts are supplied on the
  requested field, coupling and remainder parameter domains. In the
  prefix parameterization, the maps taking the current parameters to
  every earlier input's parameters are holomorphic and have images
  in those inputs' declared domains. For shared prefix coordinates
  these maps are the coordinate projections. The lists and all
  discrete choices are held fixed throughout each such domain. In the
  analytic integral presentation the sharp test is a fixed Borel
  restriction in the real free variables. The Gaussian matrices and
  the potentials admit the enlarged decoupling contours of the
  localization. Their local guards, nonvanishing logarithm domains and
  simultaneous core inclusions are part of the input. At complete
  physical data the localized sum, together with the prescribed
  scalar $\mathfrak v_{k+1}$, gives \eqref{em:eq-actual-step}.
  At an independent-current pair it is the prescribed local analytic
  expression.
\item One prescribed convergent localization gives the pointed output
  $E_{k+1}$ after the specified scalar has been separated from
  \eqref{em:eq-actual-step}. Each localized summand has a
  determining read and a fixed logarithm branch. The supplied
  localization preserves the pair gauge identities. On physical
  data these use the admitted gauge change of variables for the
  actual fibre measure, relative gauge weight and cut, including
  the prescribed chart identity if an axial presentation is used.
  Invariance of an arbitrary gauge-fixing weight is not assumed.
  Identifications of
  complete reads transport the formulas, their Borel cuts, branches
  and their assignment to carriers and points. The one-step pointing
  identity transports the complete pointed sum, allowing reassignment
  to carriers under a unit isometry that moves the
  localization partition. For an admitted isometry $r$, the real free
  coordinates change by an orthogonal map $O_{r,W}$ with
  $V'(rW,O_{r,W} B)=rV'(W,B)$. The real map may depend on $W$.
  This change transports the cut, completed
  response, measure factors and point $z$ to $rz$. There is an open
  real fluctuation ball on which all inherited backgrounds have their
  required covariance. The covariance radii $e_j$ obey
  $e_{j'}\le L^{2(j'-j)}e_j$ for $j\le j'$. This inequality compares
  thresholds for the same curvature in different units. The earlier
  backgrounds used by the actual step also have to lie in their own
  covariance classes throughout the retained real cut. The creation
  radius is $c_e\varepsilon(\theta_k)$
  at step $k\to k+1$, where $\varepsilon$ is the supplied real curvature
  threshold and $c_e$ is fixed by the domain margin. At every formed
  level, the supplied background, pair core and companion satisfy
  the domain, response, summability and cofinal compatibility
  hypotheses of
  \cref{em:lem-physical-radius,em:prop-physical-polarization,em:prop-companion-coefficient}
  uniformly on the output domains. The actual gauge and pointing
  identities supply the symmetry hypotheses as explained below.
  The readout results then define the new scalar coefficient
  before this output is used in a later list. Their analytic
  bounds use the supplied background and family estimates,
  independently of the present iteration.
\item For a list of regular entries of norm at most $E_0$, let $\mathcal
  N_k$ be its sum of the restricted original regular-family norms,
  before vacuum and quadratic subtraction, with weights
  $\omega_{k+1-j}$, plus the aged remainder norms with weights
  $\max_\iota(\bar a_{k,\iota}/\check a_{j,\iota})^2$ in the running
  chain. The weights $\omega_n$ are those in \eqref{em:eq-regular-age}.
  The comparison evaluates each regular member with its own
  counterterm. The fixed subtraction-operator bounds are included
  in $K$ and in the admission threshold, rather than being treated
  as contractive. Remainders receive their radius weights without
  that age gain.
  For $s(1+\mathcal N_k)\le\epsilon_{\rm step}$, the output is
  \begin{equation}\label{em:eq-rs-formation}
    E_{k+1}=N_{k+1}+G_{k+1,\tau_k}+F_{k+1}[\mathfrak f],\qquad
    \|F_{k+1}[\mathfrak f]\|\le
        K\{\delta(\theta_k)+s\mathcal N_k
                    +e^{-c\theta_k\delta(\theta_k)^2}\} .
  \end{equation}
  The norms are the output norms \eqref{em:eq-pointed-norm} with the
  stated radii. Fixed enlargement factors are absorbed in $K$ and in
  the threshold for $s$. Thus $s$ is $\delta(\theta_k)/a_{*,0}$ for
  the absolute chain and $\delta(\theta_k)/\bar a_{k,0}$ for the
  running chain. The separately given Gaussian families satisfy
  $\|N_{k+1}\|\le B_N$ and $\|G_{k+1,\tau}\|\le B_G|1-\tau|$ on the
  absolute cores. They depend only on the frozen Gaussian data and
  the weight ratio, and restrict to the same terms on the running
  cores. The normalization of the first family is fixed by
  \[
   \sum_{X,z}N_{k+1}(X,P_{k+1}(W),z)
        =\log Z_k(W)-\log Z_k(1),
  \]
  where $Z_k(W)$ is the supplied actual constrained reference
  Gaussian normalizer without a fluctuation cut. Its logarithm
  has the fixed branch from the positive real normalizer. This
  identity identifies its curvature coefficient with that of the
  physical Gaussian normalization. In the physical action-only
  splitting the total of the
  Gaussian mismatch is field independent, by the Gaussian gauge
  normalization. Its scalar coefficient is zero, although its
  pointed norm need not have a small-coupling factor.
\item For two admitted lists on the common smaller cores, the one-step
  comparison bounds the difference of their actual $F$ outputs by
  $Ks$ times the weighted load of the differences of their inputs,
  with the same age and radius weights. This is the comparison for
  the same completed fibre map, potential and cut domains and
  logarithm expansion. The admitted marked slot has the prescribed
  interpolation and contour domains, including the zero-marked list.
  It is additional to the one-list bound in
  \eqref{em:eq-rs-formation}.
\end{enumerate}

The constants in these inputs are uniform over the compatible tori,
over the declared parameter domains and over the number of older levels.
For every fixed prefix the compatible tori include arbitrarily large
periods in every direction, so that each fixed finite read and its
collar embed in all sufficiently large members of a cofinal sequence.
For the running construction, one sufficient size convention is
$P_i\ge2L^{11}M\check R_i$ at every used level $i$, where $P_i$
is the shortest torus period in level-$i$ units. Every common
lifted read, with its collar, must fit in a box shorter than half
each torus side. When forming level $j$ before $\theta_j$ is
supplied, the known stronger condition
$P_j\ge2L^{11}M\check R_{j-1}$ suffices for these numerical period
floors. Indeed $P_i=L^{j-i}P_j$ and
$\check R_i\le L^{j-1-i}\check R_{j-1}$ for $i<j$, while an
eventual extension has $\check R_j\le\check R_{j-1}$. These size
conditions do not supply the analytic read identities. The
standalone absolute-radius chain has no remainder-radius condition.
The blocking factor, localization side, geometric rules and regular
norm cap are fixed before raising the common frame floor. The
supplied estimates are uniform over the resulting frame-dependent
domains. Existence of these one-step data is the separate
small-field result. The conditions here specify the estimates used from
it, including the weighted comparison on the enlarged contours.

On the admitted real common covariance class, change $B$ to $O_{r,W}B$.
The Wilson term is invariant. The earlier pointed totals, their
counterterms and the admitted aged remainder totals transform as
specified. The tents satisfy $h_{rz}(rx)=h_z(x)$. By the one-step
pointing identity the transformed rooted expression therefore gives the
same pointed value at $rz$ as before at $z$. A covariance proof for the
unpointed integral alone would not give this identity. On a connected
complex parameter domain a covariance identity of holomorphic terms
continues from an open full-dimensional real slice where both are
admitted. For logarithms with already identified exponentials on that
domain, a single real anchor fixes the difference in $2\pi i\mathbb Z$.
This is also the branch check for an inherited remainder whose supplied
exponential is invariant. It does not form the remainder.

\emneedspace{12\baselineskip}
\begin{proposition}[Restriction of the same step]\label{em:prop-common-restriction}
Assume the one-step inputs on the absolute and running cores just
specified. On every complete common read the unmarked construction on
the absolute core restricts to the construction made with its restricted
unmarked list on the running core. The full and restricted unmarked
outputs admit the marked difference
\begin{equation}\label{em:eq-marked-difference}
 E_{k+1}^{\mathrm m}=E_{k+1}-E_{k+1}^{\mathrm u}|_{\bar{\mathbf a}_k}.
\end{equation}
The two Gaussian families in \eqref{em:eq-rs-formation} cancel in this
difference. At input level $j$ its regular marked member carries the
difference of the two coefficients, and its nonregular inputs are
exactly the admitted aged remainders. These identities hold for the
companions as well as the torus families.
\end{proposition}

\begin{proof}
First compare the unmarked step on the two cores. The initial list is
empty. At a later step restrict each earlier unmarked function and its
vacuum value by the previous identification. Its coefficient is a
readout of the same companion on the common near-zero physical ball.
Hence its
Wilson counterterm is unchanged. On a common read the actual response
and the arguments of these functions agree on the whole fluctuation
cut and the prescribed decoupling contours. The same term formula and
branch therefore give the same integrand in the real free coordinates.
Their integrals and the normally convergent localized series agree. This
proves the restriction at the next level. It compares analytic functions
on their cores directly. Agreement of totals on the graph of real
backgrounds would not identify functions of an independent current.

In the full versus unmarked comparison the two lists give different
potentials but have the same reference measure and cuts. Put
$D=S_{\mathfrak f}-S_{\mathfrak f^{\rm u}}$ on their common reads.
Move the last two terms of \eqref{em:eq-actual-step} inside the
exponential. The marked vertex is then exactly
\[
             D(U_k(V'(W,B)))-D(U_{k+1}(W)).
\]
In the localized logarithm put this carried difference in its marked
slot and set the slot to zero. The subtracted background value is
included once. This is the difference to which the one-step comparison
applies; a raw insertion of $D(U_k(V'))$ alone would give a different
functional. The prescribed
expansion at zero is precisely the restricted unmarked one. At this
zero value the two unmarked presentations on physical pairs give the
same nonvanishing integral, with the same real logarithm at the
reference. Any difference of their two continuations of this logarithm
is locally constant with values in $2\pi i\mathbb Z$. It vanishes on
the connected component containing the reference. On any other
admitted common component, agreement uses the matching branch
specified by the transported local formulas in the one-step input. The full output uses the prescribed branch at the full
marked input. For local independent-current terms the comparison uses
the same analytic coefficient expressions and branch choices from the
input. This gives \eqref{em:eq-marked-difference} and cancels $N$ and
$G_{\tau_k}$. Linearity of the coefficient and vacuum subtraction gives
the list difference with the marked members' own counterterms and all
and only the aged remainder inputs.

The same argument compares two tori on a common lifted read. All
coefficients transported by that lift use identical data and indices.
Choose any compatible torus large enough to contain the lifted read
to define the companion there. The comparisons prove independence of
that choice. Larger tori exhaust the finite reads. This also proves
compatibility with restriction of a frame prefix. Wrapping terms retain
their torus formulas and do not define extra companion coefficients.
\end{proof}

\emneedspace{22\baselineskip}
\begin{proposition}[Formation along prescribed prefixes]
\label{em:prop-prefix-formation}
Assume the preceding one-step inputs uniformly on the prescribed
parameter domains. Suppose the regular age weights have finite sum
$\Omega$ and the aged remainder loads of the running prefixes
are bounded by $R_*$. Fix $E_0>B_N$ before increasing the common frame
floor. Suppose $\delta(\theta)\to0$, $s\to0$ in both systems, and the
failure term in \eqref{em:eq-rs-formation} is $O(\delta(\theta))$ as the
floor increases. Take a fixed relative coupling radius $0<\sigma_*\le1/2$
with $B_G\sigma_*\le(E_0-B_N)/4$ and with the one-step admission on
$|x_k/\theta_k-1|<\sigma_*$. The running parameters lie in the stated
smaller domains. Then a sufficiently high floor, uniform in the prefix
length, gives both regular chains from their empty initial lists with
norm at most $E_0$ at every creation. The unmarked output satisfies
\[
            \|E_{k+1}^{\mathrm u}-N_{k+1}-G_{k+1,\tau_k}\|
                      \le C\delta(\theta_k)
\]
on its absolute core. The outputs are holomorphic on the admitted
parameter domains, causal in the supplied frame and weights, and
compatible with the stated torus, companion and core restrictions.
The unmarked conclusions also hold on their own under the absolute
one-step hypotheses, summable regular weights and the stated tail
bound. They require neither a running chain nor inherited remainders
nor a lower drift bound; their floor depends only on these absolute
inputs.
\end{proposition}

\begin{proof}
Put $D_* =\Omega E_0+R_*$ for the two-chain assertion.
For the unmarked assertion alone put $D_*=\Omega E_0$.
Choose the floor so that for the relevant
systems
\[
    s(1+D_*)\le\epsilon_{\rm step},\qquad
    K\{\delta(\theta)+sD_*+e^{-c\theta\delta(\theta)^2}\}
                               \le (E_0-B_N)/2 .
\]
These are uniform scalar conditions on the half-line above that floor.
All earlier regular creation norms are at most $E_0$ at the induction
step. Their inherited restrictions have no larger norm. Their total
load is at most $\Omega E_0$, and the full list adds at most
$R_*$. The one-step hypotheses thus apply to the actual lists.
Equation \eqref{em:eq-rs-formation} gives the new norm at most
$B_N+(E_0-B_N)/4+(E_0-B_N)/2<E_0$. The empty input at the first step
starts the same induction. With no remainder terms and fixed absolute
radii, $s=\delta/a_{*,0}$ and the bound for $F$ is $C\delta$ as claimed.
For example the failure condition follows when
$\delta(\theta)=c_\delta\theta^{-1/2}(\log\theta)^{\mathfrak p_0}$
with $c_\delta>0$ and $2\mathfrak p_0>1$. Its negative exponent then
dominates any fixed multiple of $\log\theta$. On running radii the
required smallness of $s$ follows from their stated comparison with
$\delta$, independently of this failure estimate.

At step $k$ the parameters of the earlier outputs enter only through
their already formed analytic functions. The supplied parameter maps
land inside those earlier domains, so their compositions are defined
and holomorphic. The fixed-domain step and
its normally convergent expansion give joint holomorphy with the new
parameters. Finite induction gives it on each declared prefix domain.
There is no sum over all possible prefixes. The new family, its
coefficient and its creation covariance radius use only data through
$\theta_k,x_k$. The choice of the next frame entry is needed only to
specify the inherited restriction. The common read identities in
\cref{em:prop-common-restriction} persist at each induction step. For
the absolute chain no lower drift bound and no remainder are used.
Its independent weight parameters may range over the full admitted
product of relative discs. The marked parameter domain projects into
that product in the shared weights. For example a tube
$|z_i|<\varpi_0\theta_i(\log\theta_i)^{-2q_0}$ with
$q_0\ge0$, $\theta_i>e$ and $0<\varpi_0\le\sigma_*/2$ has this
inclusion. Its inherited functions retain their separate declared
domains. For the full chain the prescribed weaker
flow and smaller analytic domains are used only where required for
the running radii and inherited loads.
\end{proof}

\section{The marked estimate}
\label{em:sec-marked-recursion}

Write $u_i$ for the norm of the marked difference of the two
outputs at level $i$, on their common running creation domain,
including the prescribed parameter supremum. The level-$i$ output was formed using $\theta_{i-1}$.
In later lists it is used by restriction and with its own linear
quadratic subtraction. The restriction and subtraction norms used
by the one-step comparison are included in its fixed constants.

\emneedspace{18\baselineskip}
\begin{proposition}[Uniform marked recurrence]\label{em:prop-marked-recursion}
Use \eqref{em:eq-regular-age}--\eqref{em:eq-running-weight} and
the inherited loads of \cref{em:lem-remainder-load}. Suppose the
two constructions start with the empty list and satisfy the
one-step comparison and common-restriction hypotheses of
\cref{em:prop-common-restriction}. At every step $m+1$, $0\le m\le N$,
they satisfy
\begin{equation}\label{em:eq-marked-input}
  u_{m+1}\le K_M\mathsf s_m\left(
         \sum_{i=1}^m\omega_{m+1-i}u_i
                   +\sum_{j\in J_m}w_{m,j}\rho_j\right).
\end{equation}
Here $K_M$ is uniform on the stated one-step admission class and
is fixed before choosing the frame floor. Above a sufficiently
large fixed floor $X$, one has
\begin{equation}\label{em:eq-marked-envelope}
  \sup_{1\le i\le N+1}\theta_{i-1}^{p}u_i
        \le 2K_MC_{\mathfrak h}\mathsf s_* ,\qquad
  u_{m+1}\le 2K_MC_{\mathfrak h}\mathsf s_m\theta_m^{-p},
  \qquad
  \mathsf s_* =\sup_{y\ge X}
    \frac{\delta(y)}{(1+\beta_{\rm s})\alpha_0(y)}.
\end{equation}
If the norms and input estimate hold uniformly on an admitted
parameter domain, the bounds are uniform on that same domain.
\end{proposition}

\begin{proof}
By \cref{em:prop-common-restriction}, the two Gaussian families
cancel. Linearity of the coefficient and of vacuum subtraction
makes the difference of the carried regular actions the action
of the marked regular members with their own counterterms. The
remaining input differences are the admitted aged remainders.
The one-step comparison charges their original family norms,
with the fixed subtraction costs included in $K_M$. It gives
\eqref{em:eq-marked-input} on the same actual cores.

Let $u_m^*=\max_{1\le i\le m}\theta_{i-1}^pu_i$, with $u_0^*=0$.
For every $i\le m$, the lower flow comparison and the floor give
$\theta_{i-1}\ge\theta_m/2$. Hence
\[
   \theta_m^p\sum_{i=1}^m\omega_{m+1-i}u_i
                        \le 2^p\Omega u_m^*.
\]
Multiplying \eqref{em:eq-marked-input} by $\theta_m^p$ and using
\eqref{em:eq-aged-load} gives
\[
  \theta_m^p u_{m+1}
       \le K_M\mathsf s_m(2^p\Omega u_m^*+C_{\mathfrak h}).
\]
The ratio defining $\mathsf s_*$ is a fixed multiple of
$(\log y)^{\mathfrak p_0-q}$. Thus we may choose $X$ so that
$K_M\mathsf s_*2^p\Omega\le1/2$. Inductively, if
$u_m^*\le 2K_MC_{\mathfrak h}\mathsf s_*$, the right side is at
most $2K_MC_{\mathfrak h}\mathsf s_m$. This proves both
estimates at the new level and starts from $u_0^*=0$.
\end{proof}

Only frame data through $m$ are used in the estimate of the newly
formed output. The later restriction involving $\theta_{m+1}$ is
performed on that output once this entry is available. The recent
load \eqref{em:eq-recent-load} bounds the available input sizes in
that band. Its use in forming a new remainder or a boundary family
requires the separate operations for those outputs.

\section{Extraction on the physical radii}
\label{em:sec-coefficient}

The coefficient is a second derivative in the physical block datum.
In level-$j$ units write $\eta_j\in(0,1]$ for the fine mesh, with
$\eta_j^{-1}$ integral, and
$V=e^{iB}$ for the unit-lattice block datum. The stored positive
bonds specify all signs of the cochain coboundary. In a chart based
at the identity pair $P_\circ=(\mathbf 1,0)$ use
\begin{equation}\label{em:eq-physical-pair-norm}
 P=(e^{i\eta_j A},J),\qquad
 \|(A,J)\|_{1,j;D}
   =\max\{\|A\|_{\infty,D},\|\nabla_{\eta_j}A\|_{\infty,D},
                             \|J\|_{\infty,D}\} .
\end{equation}
The differences in this formula are physical difference quotients
on the complete stencils of the declared read window $D$. The current
$J$ is independent in the complex pair domain. On a physical
background it is the Wilson current in the physical fine pairing.
The two radii $r_0,r_1$ bound the physical curvature and current, and
the $C^1$ perturbation of the potential, respectively.

\begin{lemma}[A physical ball in the pair domain]
\label{em:lem-physical-radius}
Suppose the hereditary pair domain at radii $r_0,r_1$ has the following
tower input. For a triple whose real and complex plaquette defects
are at most $\eta_j^2\epsilon_0$, whose real factor admits the declared
local gauges, whose perturbation has physical $C^1$ norm at most
$\epsilon_1$ and whose independent current is at most $\epsilon_0$,
the prescribed real and complex towers on every allowed restricted
window exist and satisfy the hereditary curvature and physical-current
bounds with radius $K_T(\epsilon_0+\epsilon_1^2)$ in the respective
units. Assume this estimate with uniform constants and on fixed
positive ceilings for $\epsilon_0,\epsilon_1$.
Then the ball
\[
       \|(A,J)\|_{1,j;D}< a,
       \qquad a=c\min\{r_0,r_1,1\}
\]
about $P_\circ$ lies in the pair domain on $D$. Here $c>0$ can be
chosen independently of the fine mesh, level and compatible volume.

Suppose also that the actual continued physical background pair
$p_j(B)$ is holomorphic on $\|B\|_\infty<r_B$, satisfies
$p_j(0)=P_\circ$, and has $\|p_j(B)\|_{1,j;D}\le C_B\|B\|_\infty$
in the coordinates of \eqref{em:eq-physical-pair-norm}, with uniform
$r_B,C_B>0$. Its data ball of radius
$\min\{r_B/2,a/(2C_B)\}$ then lies in the same pair domain and
is comparable to $a$ with fixed constants.
\end{lemma}

\begin{proof}
Use $\mathbf 1$ as the real factor. It has zero curvature and the
zero potential in each of the required local gauges. In the product
of four link exponentials around a plaquette, the linear term is
$i\eta_j^2\mathrm d_{\eta_j}A$. Here $\mathrm d_{\eta_j}$ is the
oriented curl per unit length. The remaining series is bounded by
$C\eta_j^2\|A\|_\infty^2$ on the fixed small chart, because every
monomial of order at least two contains two factors of size at most
$\eta_j\|A\|_\infty$. Thus the plaquette defect is bounded by
$C\eta_j^2(a+a^2)$. The independent current is bounded by $a$.
Use $\epsilon_0=C'(a+a^2)$ and $\epsilon_1=a$ in the tower estimate.
Reducing $c$ places these numbers below the input ceilings and makes
$K_T(\epsilon_0+\epsilon_1^2)<r_0$. It also ensures the simple pair
conditions at $r_0,r_1$. This verifies every hereditary condition
on its own window. The assertion about $p_j$ follows by the displayed
linear bound on the actual background. The first difference in
\eqref{em:eq-physical-pair-norm} is used in the plaquette estimate.
\end{proof}

In the physical chart write $H_j$ for the first derivative of the
potential of $p_j$. The first response of the pair is
\begin{equation}\label{em:eq-coefficient-response}
       \mathfrak h_j B
       = (H_j B,\mathrm d_{\eta_j}^*\mathrm d_{\eta_j}H_j B).
\end{equation}
Indeed the Wilson action to second order is
$\tfrac12\eta_j^4\sum_p|(\mathrm d_{\eta_j}A)(p)|^2$ in four
dimensions. Its physical gradient uses the bond pairing with weight
$\eta_j^4$. Thus the second entry in
\eqref{em:eq-coefficient-response} is the derivative of the current
on link fields. All these normalizations use the half-trace colour
pairing.

\subsection{Polarization of a pointed family}

Use localization domains $X$ of fixed cell side $M$ and denote their
tree length in cell units by $d_j(X)$. A base point $z\in X$ is a
unit-lattice site. Write $X^+$ for the fixed-width enlarged read window
of $X$. We use the rooted entropy bound $C e^{\kappa_D m}$ for the
number of domains meeting a given cell with $m\le d_j(X)<m+1$.
The diameter in level-$j$ units is at most $C M(m+1)$ and the number
of base points is at most $C_M(m+1)^4$. These hypotheses hold on the
declared tori and on the companion, with the distance and lift used
in the family.

For a pointed family $E(X,P,z)$, let $\mathcal A$ be a bound for its
weighted supremum norm with exponent $\kappa$. Thus each term on
its core has modulus at most $\mathcal A e^{-\kappa d_j(X)}$.
The family is holomorphic in the pair. Use an orthonormal colour
basis $t_1,t_2,t_3$,
and write $\delta_b^s$ for the unit datum on the positive bond $b$
in colour $t_s$. For the companion, define
\begin{equation}\label{em:eq-polarization-kernel}
 \Pi_E(b,c)=\frac13\sum_{s=1}^3\sum_X\sum_{z\in X}
       D_P^2 E(X,P_\circ,z)
             [\mathfrak h_j\delta_b^s,\mathfrak h_j\delta_c^s] .
\end{equation}
The same definition is used on a torus, with its own physical
response. The Hessian is in the coordinates
\eqref{em:eq-physical-pair-norm}.

\emneedspace{22\baselineskip}
\begin{proposition}[The radius cost in the polarization]
\label{em:prop-physical-polarization}
Let $\kappa>\kappa_D$ with a fixed margin. Suppose the pair ball of
radius $a>0$ in \eqref{em:eq-physical-pair-norm} on each read window
lies in the core, and the actual first response satisfies
\begin{equation}\label{em:eq-response-decay}
 \|\mathfrak h_j\delta_b^s\|_{1,j;X^+}
                \le C_h e^{-\delta\operatorname{dist}(b,X^+)}
\end{equation}
for some fixed $\delta>0$. The kernel
\eqref{em:eq-polarization-kernel} converges absolutely and obeys
\begin{equation}\label{em:eq-polarization-bound}
  |\Pi_E(b,c)|\le C a^{-2}\mathcal A
                            e^{-\delta_1|b_--c_-|} .
\end{equation}
Here the distance is periodic for a torus. At a single base point
$z$, the absolute double sum of the corresponding kernel with weight
$|b_--z|\,|c_--z|$ is bounded by $Ca^{-2}\mathcal A$.
The same bound, with a constant depending on the powers, holds
with any fixed nonnegative powers of $1+|b_--z|$ and $1+|c_--z|$.
The constants depend on the fixed geometry and response and entropy
bounds and are uniform in depth and compatible volume. On a finite
torus, if the complete pair total is invariant under constant colour
conjugations near $P_\circ$, then the kernel is the colour average
of the Hessian of the composed physical total.
\end{proposition}

\begin{proof}
For two nonzero directions $v,w$ in a read window, take the bidisc
with radii $a/(4\|v\|_{1,j;X^+})$ and
$a/(4\|w\|_{1,j;X^+})$ in the two direction parameters. Its image
lies in the pair ball. The two-variable Cauchy formula yields
\[
 |D_P^2E(X,P_\circ,z)[v,w]|
   \le 16a^{-2}\mathcal A e^{-\kappa d_j(X)}
                  \|v\|_{1,j;X^+}\|w\|_{1,j;X^+} .
\]
The estimate is immediate if either direction is zero. On the
response directions the two factors contribute
$e^{-\delta\operatorname{dist}(b,X^+)}$ and
$e^{-\delta\operatorname{dist}(c,X^+)}$. The distance between the
lower endpoints is at most the sum of these distances and
$C_M(d_j(X)+1)$. Reserve a fixed small fraction of all three
exponentials to supply $e^{-\delta_1|b_--c_-|}$. Choose the fraction
so that the remaining tree exponent is larger than $\kappa_D$.
Root each domain at a cell nearest $b$. At a fixed root, the domain
entropy and the base point count are summable against the remaining
tree exponent. The remaining root-distance exponential sums over
root cells, also on a torus by counting shortest lifts. This proves
\eqref{em:eq-polarization-bound}.

For the single base point bound, root at the cell containing $z$.
For a fixed $X$, group $b$ by a cell of $X^+$ nearest its lower
endpoint. The weight $|b_--z|$ is at most its distance to that cell
plus $C_M(d_j(X)+1)$. The cell count in $X^+$ and the weighted
distance sum are bounded by polynomials in $d_j(X)+1$. Use the same
bound for $c$. The tree exponent exceeds the entropy exponent, so
the resulting polynomial factors sum. This proof makes no symmetry
assertion about the kernel at a single base point.

For the finite-torus identification, the differential of the pair
total at $P_\circ$ is invariant under the adjoint action of
$\operatorname{SU}(2)$. Differentiating this invariance gives zero on
every commutator in each potential and current coordinate. Since
commutators span the colour algebra, the differential vanishes.
The chain rule for the composed total therefore has zero contribution
from the second derivative of the physical background. The remaining
term is precisely \eqref{em:eq-polarization-kernel}. This argument
uses only the colour symmetry of the complete pair total.
\end{proof}

If the families vary holomorphically with admitted parameters and
have the same uniform estimates on compact parameter sets, the
derivatives and the sums in \eqref{em:eq-polarization-kernel} are
holomorphic. This follows by the same Cauchy contours and normal
convergence. For a given frame the response is fixed. The conclusion
also applies to a holomorphic response with the same local bounds.

\emneedspace{27\baselineskip}
\subsection{The scalar coefficient and its normalization}

\begin{proposition}[The companion coefficient]
\label{em:prop-companion-coefficient}
Assume the hypotheses of \cref{em:prop-physical-polarization} for
the finite families and their companion. On each common nonwrapping
lift the actual functions and their pair charts agree with the
companion, including their read collars. Assume a cofinal sequence
of compatible cubic tori whose periods tend to infinity, and local
convergence of the responses on every
bounded window along it, with the uniform bound
\eqref{em:eq-response-decay}. Suppose the finite pair totals are
invariant under constant colour conjugations as above, and the
full physical totals on
these cubic tori are invariant under the unit translations,
signed coordinate permutations and physical gauge transformations
near the trivial datum. Then the companion kernel is translation
invariant, written $\Pi_{E,\mu\nu}(u)$ with $u=b_--c_-$ and bond
directions $\mu,\nu$. Its zeroth and first spatial moments vanish.
For $\mu\ne\nu$ set
\begin{equation}\label{em:eq-coefficient-moment}
    b_j(E)=\sum_{u\in\mathbb Z^4}\Pi_{E,\mu\nu}(u)u_\mu u_\nu .
\end{equation}
This value is independent of $\mu,\nu$, and
\begin{equation}\label{em:eq-maxwell-moment}
 \sum_u\Pi_{E,\mu\nu}(u)u_\rho u_\sigma
  =b_j(E)(\delta_{\mu\rho}\delta_{\nu\sigma}
      +\delta_{\mu\sigma}\delta_{\nu\rho}
      -2\delta_{\mu\nu}\delta_{\rho\sigma}) .
\end{equation}
In particular $|b_j(E)|\le C a^{-2}\mathcal A$.
The map $b_j$ is linear and real on real data. It is unchanged by
admitted restriction of the same functions around $P_\circ$ and by
changing coherent localizations of the same physical total.
\end{proposition}

\begin{proof}
Fix two bond arguments and first retain domains of bounded tree
length whose nearest root to the first argument lies in a bounded
set. Eventually each such window and collar embeds in the torus.
Its Hessian at $P_\circ$ agrees with the companion Hessian by the
common-chart hypothesis. The responses converge there. The tails
of the rooted sums are uniform by the proof of
\eqref{em:eq-polarization-bound}. Sending both bounds to infinity
therefore identifies the infinite kernel with the limit of the
physical torus Hessian kernels. Their spatial moments through order
two pass to this limit by the exponential bound, using the short
representatives on the torus. These steps also show that two
coherent localizations of the same total give the same kernel.
Restriction of an actual function to a smaller neighborhood of the
origin preserves its derivatives.

The covariance identities of the finite physical Hessians pass to
the limit. The gauge identity implies that the Hessian annihilates
a gradient in either slot. In fact, differentiate the physical
gauge identity first in a site gauge and then in a bond datum.
At zero the gauge tangent is the ordinary signed site difference,
and the derivative of that tangent is multiplied by the zero first
differential. The statement passes to the infinite kernel on
finitely supported site functions. Colour averaging preserves it.

Write $\widehat\Pi(p)=\sum_u e^{-ip\cdot u}\Pi_E(u)$ and
$t_\mu(p)=e^{ip_\mu}-1$. The exponential bound gives derivatives of
the symbol near zero, and the Ward identity reads
$\widehat\Pi(p)t(p)=0$. Its term of degree one gives
$\widehat\Pi(0)p=0$ for every $p$, hence $\widehat\Pi(0)=0$.
Put $U_{\mu\mu}(p)=e^{ip_\mu/2}$ and use the midpoint symbol
$U(p)^{-1}\widehat\Pi(p)U(p)$. In these coordinates every signed
coordinate permutation acts by its ordinary signed matrix on the
link directions, since it maps bond midpoints to bond midpoints.
Inversion therefore makes this symbol even. This eliminates its
linear term, and the changes of phase preserve the degree-two
part because the constant and linear parts vanish.

Let $Q(p)$ be that quadratic part. Coordinate reflections force
$Q_{\mu\mu}=A p_\mu^2+B\sum_{\nu\ne\mu}p_\nu^2$ and
$Q_{\mu\nu}=C p_\mu p_\nu$ for $\mu\ne\nu$. Permutations give
common coefficients $A,B,C$. The leading midpoint Ward vector is
$ip$, so $Q(p)p=0$. Its $\mu$ component gives $A=0$ and $C=-B$.
Thus $Q(p)=B(|p|^2 I-pp^{\mathsf T})$. The second derivative of
the Fourier series is the negative second moment. The mixed
off-diagonal moment in \eqref{em:eq-coefficient-moment} is $B$,
and the full tensor is \eqref{em:eq-maxwell-moment}. The coefficient
bound follows by summing $C a^{-2}\mathcal A |u|^2e^{-\delta_1|u|}$.
Linearity and reality follow from the defining derivatives.
\end{proof}

The limiting response may be taken in any coherent representative
family for which the hypotheses of
\cref{em:prop-companion-coefficient} hold. Suppose two such
responses in the common pair chart differ, on each datum direction
\(B\), by \((D_{\eta_j}\lambda_B,0)\), where the restriction
of \(\lambda_B\) to each determining window generates an admitted
fine gauge there. For this site-to-link difference,
\(D_{\eta_j}\lambda(x,x+\eta_je_\mu)
=(\lambda(x)-\lambda(x+\eta_je_\mu))/\eta_j\). At \(P_\circ\) the
gauge tangent in the current slot is zero, because the current
transforms by conjugation and \(J=0\). Differentiate the gauge
identity of a local pair term first in the gauge and then in an
arbitrary pair direction. Its first differential at \(P_\circ\)
vanishes by constant-colour invariance. The resulting Hessian
therefore annihilates \((D_{\eta_j}\lambda_B,0)\) in either
slot. Thus both responses give the same term in
\eqref{em:eq-polarization-kernel}, and hence the same coefficient.
This comparison uses actual response convergence for both
representative families.

For the unit-lattice Wilson action the quadratic term is
$\tfrac12\|\mathrm d B\|_2^2$. Its Hessian symbol is
$|t|^2 I-tt^*$, whose quadratic part is the Maxwell symbol with
coefficient one. This fixes $b_j(A_{\rm block})=1$, with no additional
factor of two. For a fine Wilson action pulled back by a background,
the same conclusion has a separate hypothesis and warrant.
Assume the curl response $\mathrm d_{\eta_j}H_j$ is covariant under
unit translations and $H_j$ reproduces affine curls exactly in physical units. That is,
$\mathrm d_{\eta_j}H_jB$ is the constant continuum curl of an affine
unit-lattice datum $B$. Interpret these affine data by the absolutely
convergent response moments. The curl of the response to a midpoint
plane wave in direction $\mu$ then has first-order symbol
$i p\wedge e_\mu$ at every fine position in a unit cell. To see this,
expand the plane wave at the cell origin. Its constant part has zero
curl, its affine part has the stated curl by reproduction, and the
second-order remainder is bounded by the second absolute moment of
the response. That moment is uniform by
\eqref{em:eq-response-decay}. Squaring the first-order term in the
fine energy and summing with weight $\eta_j^4$ over the fine
positions in a unit cell gives $|p|^2 I-pp^{\mathsf T}$, since the
fine weights sum to one for each plaquette orientation. Thus
the pulled-back fine Wilson action has coefficient one under this
affine response hypothesis.

For the prescribed pointed Wilson insertion the tents sum to one
on the physical total. Each tent has bounded support in level units.
On a pair core of curvature radius $r_0$, a fine plaquette action is
$O(\eta_j^4 r_0^2)$ by the trace expansion at the identity. The
$O(\eta_j^{-4})$ fine plaquettes in a tent give a pointed norm at
most $C r_0^2$, with the fixed localization weight included in $C$.
When $r_0$ and $a$ are comparable and the insertion has coefficient
one by the preceding normalization, subtracting
$b_j(E)$ times this insertion costs at most $C\mathcal A$ in the
same pointed norm. The scalar normalization uses the total and
does not identify the tensor at each base point.

For $g>0$, if the unit-block datum is written $B=gB_{\rm s}$, the chain rule gives
$D^2(E\circ p_j(g\,\cdot))(0)=g^2D^2(E\circ p_j)(0)$. The scaled
coefficient is $g^2b_j(E)$, while the Cauchy radius in the scaled
datum is proportional to $a/g$. Its bound is thus
$|g^2b_j(E)|\le C\mathcal A(g/a)^2$. Dividing by $g^2$ gives the
same physical estimate $C a^{-2}\mathcal A$. The bound on the
physical field cannot be replaced by a bound on the normalized
fluctuation variable.

\subsection{The marked coefficient on a prescribed prefix}

At step $k\to k+1$, use the common absolute and running restrictions
of the two actual outputs. On the absolute domain the unmarked
output decomposes as $N_{k+1}+G_{k+1,\tau_k}+F_{k+1}^{\rm u}$.
Here $N_{k+1}$ is the normalized Gaussian family, and
$G_{k+1,\tau_k}$ is the family of the change of Gaussian weight.
Throughout the admitted physical chart its total is a background-free
number by the action--gauge splitting in
\cite[Proposition~4.3]{GI}. Consequently its coefficient
vanishes by \cref{em:prop-companion-coefficient}. This cancellation
uses the splitting identity for the actual Gaussian and does not
follow from the size of $G_{k+1,\tau_k}$.

\emneedspace{25\baselineskip}
\begin{theorem}[The coefficient bounds on the two domains]
\label{em:thm-coefficient-envelope}
Use the pair-domain, response and companion hypotheses of
\cref{em:lem-physical-radius,em:prop-companion-coefficient} uniformly
on the corresponding formed families. Suppose the unmarked
residual has norm at most
\[
  C_{\rm u}\theta_k^{-1/2}(\log\theta_k)^{\mathfrak p_0}
\]
on its absolute pair core, holomorphically in the independent
weight deviations $z_i=\theta_i-x_i$ on
$|z_i|<\sigma_*\theta_i$, $0\le i\le k$. The normalized Gaussian
family is fixed on the frame. Suppose the actual marked difference
$E_{k+1}^{\rm m}$ in \eqref{em:eq-marked-difference}
has norm at most $C_{\rm m}\theta_k^{-p}$ on the common running
core, where $p=(\kappa_0-2)/2$ and $\kappa_0>4$. Assume joint
holomorphy on the admitted narrower weight tube and inherited
remainder domains. The weight tube is contained in
$|z_i|<\varpi_0\theta_i(\log\theta_i)^{-2q_0}$ for fixed positive
$\varpi_0,q_0$. Let the frame entries exceed $e$. The two
component radii for that creation step are comparable to
$\theta_k^{-1/2}(\log\theta_k)^q$ with fixed positive comparison
constants. The absolute radii are fixed positive numbers.
Then
\begin{equation}\label{em:eq-coefficient-envelope}
 \begin{aligned}
 b_{k+1}&=b^0_{k+1}+f_{k+1}(z_0,\ldots,z_k)+q_{k+1},\\
 |f_{k+1}|&\le C\theta_k^{-1/2}(\log\theta_k)^{\mathfrak p_0},\\
 |q_{k+1}|&\le C\theta_k^{-(\kappa_0-4)/2}(\log\theta_k)^{-2q}.
 \end{aligned}
\end{equation}
The last estimate and its holomorphy hold on the admitted marked
parameter domain. The coefficients are real on real data. They
are independent of the compatible torus and of subsequent
restrictions of the same creation functions. If
\eqref{em:eq-marked-envelope} holds on this domain, the bound
for $q_{k+1}$ also holds with the extra factor $\mathsf s_k$
defined in \eqref{em:eq-running-weight}.
\end{theorem}

\begin{proof}
Set $b_{k+1}^0=b_{k+1}(N_{k+1})$,
$f_{k+1}=b_{k+1}(F_{k+1}^{\rm u})$ and
$q_{k+1}=b_{k+1}(E_{k+1}^{\rm m})$. Linearity and the vanishing
coefficient of $G_{k+1,\tau_k}$ give the equality. The absolute
core has a fixed physical ball by \cref{em:lem-physical-radius},
so the extraction bound costs a fixed constant for $f_{k+1}$.
The running core has physical ball radius at least
$c\theta_k^{-1/2}(\log\theta_k)^q$ in the same lemma. Its
extraction cost is at most $C\theta_k(\log\theta_k)^{-2q}$.
Multiplication by $\theta_k^{-p}$ gives the stated marked bound.
Applying the sharper norm estimate in
\eqref{em:eq-marked-envelope} before extraction retains $\mathsf s_k$.
This spends the physical radius once. When a transferred old
quadratic term is charged with the factor $(a/a_i)^2$ on the
output radius $a$, the readout costs
$a^{-2}(a/a_i)^2=a_i^{-2}$. A separate charge of the creation
radius would count it twice. The parameter claims follow from
the normal convergence in \cref{em:prop-physical-polarization}
and the companion identification in
\cref{em:prop-companion-coefficient}. The coefficient uses
the step-$k$ creation radii and therefore requires no next frame
entry.
\end{proof}

\begin{corollary}[Cauchy bounds in the weight and bare variables]
\label{em:cor-coefficient-slopes}
Under \cref{em:thm-coefficient-envelope}, let $\epsilon_k$ denote
the unmarked bound in \eqref{em:eq-coefficient-envelope} and put
$\|v\|_\theta=\max_{0\le i\le k}|v_i|/\theta_i$.
On the half-polydisc $|z_i|\le\sigma_*\theta_i/2$ one has
\[
 |Df_{k+1}(\mathbf z)[v]|\le C_{\sigma_*}\epsilon_k\|v\|_\theta.
\]
In particular, when $\theta_i\ge\theta_k-c_5$ on the prefix and
$\theta_k\ge2c_5$, the Lipschitz bound in the unweighted supremum
norm is at most $C\epsilon_k/\theta_k$ on that half-polydisc.

For the marked term, suppose the actual fixed-frame outputs and
inherited lists are holomorphic along a bare-variable disc
$|\zeta-\zeta_0|<R$ with the indicated common domains, cuts,
branches and companion. Suppose the marked bound in
\eqref{em:eq-coefficient-envelope} holds throughout this disc.
Writing $\epsilon'_k$ for that bound, one has
\[
       |\partial_\zeta q_{k+1}(\zeta)|\le 2\epsilon'_k/R
                    \qquad (|\zeta-\zeta_0|\le R/2).
\]
\end{corollary}

\begin{proof}
For nonzero $v$ the line $\mathbf z+t v$ remains in the unmarked
polydisc for $|t|<\sigma_*/(2\|v\|_\theta)$. Cauchy's formula on
that line gives the derivative bound. The zero direction is
immediate. Integrating along a segment in the convex half-polydisc
gives the Lipschitz bound, and the lower frame comparison bounds
$\|v\|_\theta\le 2\|v\|_\infty/\theta_k$. This argument avoids a
sum of bounds for the separate coordinate derivatives.
On the marked disc use the actual one-variable holomorphic
coefficient, including the inherited dependence. The distance
from its half-disc to the outer boundary is at least $R/2$.
Cauchy's formula gives the second bound. The line in the unmarked
polydisc gives no enlargement of the marked bare-variable disc.
\end{proof}

For example, two inherited lists at real centres of the same
frame may be compared on their actual common cores. The list
$H_\upsilon=(1-\upsilon)H_0+\upsilon H_1$ has norm at most $4r$
on $|\upsilon|<3/2$ if each has norm at most $r$ there before
interpolation. The estimate is the triangle inequality. Their
common companion and linear covariance identities persist under
this interpolation. The list can be used in the marked estimate
when its slot admits that norm on those same cores. Holomorphy
along this list interpolation and along a bare-variable disc are
different hypotheses on the actual inputs.

\section{Forward iteration of complete regular profiles}
\label{em:sec-profile}

The list estimates also have a formulation for the complete regular
profile after its scalar curvature coefficient has been subtracted.
This formulation requires a transport estimate on complete functions.
An estimate on one chosen local representation alone does not imply
the required composition identity.

Fix a finite horizon $N$. For $0\le j\le N$, let $\mathcal B_j$ be a
complex Banach space of complete pointed analytic families at scale
$j$, using the complex classes just defined. Each family has zero value and first derivative at the trivial
field, zero companion curvature coefficient, and the prescribed gauge,
lattice-symmetry and period-compatibility identities. The spaces use
a common physical field radius and a common localization rate.
Their norms include the companion and finite-period compatibility
bounds. A faithful realization starts with a Banach space of weighted
summable local analytic representations, with bounded complete
evaluation and compatibility maps. Quotient this Banach space by
the joint kernel of all complete evaluations. That kernel is closed,
so the quotient is Banach and two quotient classes give the same
complete functions only when they are equal. The gauge, covariance,
zero-jet and zero-coefficient conditions are imposed as closed
identities, using their bounded maps on the quotient. The supplied
local formation construction must provide these bounded maps and
the compatible analytic representations.

\begin{lemma}[Affine test of a complete pointing]
\label{em:lem-affine-pointing}
On the infinite unit lattice let \(f_z\) be a complete pointed
physical family that is Fr\'echet holomorphic near zero in the
supremum norm on bounded bond data. Assume its physical
Hessian kernels have absolutely summable moments of every fixed
polynomial order in the distances of their two bond arguments from
\(z\), uniformly in \(z\). Define its bulk kernel by the translation
sum of these Hessians and its colour-averaged coefficient \(b(f)\)
by \eqref{em:eq-coefficient-moment}. The pair hypotheses of
\cref{em:prop-physical-polarization,em:prop-companion-coefficient}
are one sufficient source of these moment bounds.
Suppose each \(f_z\) is gauge invariant, and the family is covariant
as a scalar pointing under all unit translations and signed
coordinate permutations. Let \(F_{\rho\mu}=-F_{\mu\rho}\) be a
real constant colour two-form and \(a^F_\mu(x)=
\frac12\sum_\rho(x_\rho-z_\rho)F_{\rho\mu}\) its affine potential.
Then the affine test, defined by the absolutely convergent Hessian
moments, is
\begin{equation}\label{em:eq-affine-pointing}
 \frac12D^2 f_z(0)[a^F,a^F]
                  =\frac12 b(f)\sum_{\rho<\mu}|F_{\rho\mu}|^2 .
\end{equation}
\end{lemma}

\begin{proof}
The first differential of the complete pointing vanishes by constant
colour conjugation. Differentiate its gauge invariance in a site
gauge and in a bond datum. Its Hessian annihilates a gradient in
either slot. Constant colour conjugation makes every colour block
scalar, since the adjoint representation is the full rotation
group on the three real colours. Let \(K_z(b,c)\) be this scalar
Hessian. The assumed absolute moments allow the Ward identities to
be tested on
gradients of linear or quadratic site polynomials by cutting them
off at radius \(R\) and sending \(R\) to infinity. The gradients
of these cutoffs are bounded by \(C(1+|x-z|)\) uniformly in \(R\).

Define the pointed second moment
\[
 T_{\mu\nu,\rho\sigma}(z)
  =\sum_{\substack{b\ \mathrm{in\ direction}\ \mu\\
                   c\ \mathrm{in\ direction}\ \nu}}
      (b_--z)_\rho(c_--z)_\sigma K_z(b,c).
\]
Testing linear site polynomials gives zero row and column sums
in each direction. Testing quadratic site polynomials then makes
\(T\) antisymmetric under \(\mu\leftrightarrow\rho\) and
\(\nu\leftrightarrow\sigma\). The extra constant in a repeated-index
discrete derivative vanishes by the linear test. The same zero sums
permit using midpoints in place of the lower endpoints. Translation
covariance makes \(T\) independent of \(z\). Using the midpoints,
a reflection acts on the four indices with the ordinary four sign
factors. Thus any nonzero component must have each index occurring
an even number of times. Antisymmetry leaves only the two pairings
\(\mu=\nu,\rho=\sigma\) and \(\mu=\sigma,\nu=\rho\) with distinct
indices. Permutations and antisymmetry give
\[
 T_{\mu\nu,\rho\sigma}
       =t(\delta_{\mu\nu}\delta_{\rho\sigma}
                  -\delta_{\mu\sigma}\delta_{\nu\rho})
\]
for one scalar \(t\).

Fix the first lower endpoint at zero in the bulk second moment and
sum over \(z\). Translation covariance lets \(b_--z\) range over
all first lower endpoints in \(K_0\) and \(c_--z\) over all second
ones. The terms involving only the first or only the second endpoint
vanish by the directional zero sums. The two cross terms give
\(-T_{\mu\nu,\rho\sigma}-T_{\mu\nu,\sigma\rho}\). For
\(\rho=\mu,\sigma=\nu,\mu\ne\nu\) this is \(t\), so \(t=b(f)\).
Contracting \(T\) with \(F_{\rho\mu}/2,F_{\sigma\nu}/2\)
and multiplying by \(1/2\) proves \eqref{em:eq-affine-pointing}.
\end{proof}

\emneedspace{8\baselineskip}
\begin{proposition}[Transport after summing old anchors]
\label{em:prop-anchor-transport}
Fix \(i<k\), let \(\xi=L^{-(k-i)}\), and suppose the literal
intermediate background map \(B_\xi\) from final to old datum is
Fr\'echet holomorphic on the admitted fixed radii in the supremum
norms on bounded bond data, preserves the zero datum,
and intertwines
physical gauges and lattice isometries on the scaled old lattice.
Assume its first response sends a final affine potential of curvature
\(F\) to an old affine potential of curvature \(\xi^2F\) modulo
the gradient of a polynomially bounded site field. This identity
includes any node frame change. Put $H_\xi=DB_\xi(0)$ and assume
its first absolute row moment is finite,
\[
 \sup_b\sum_c\|H_\xi(b,c)\|
              (1+|c_--\xi b_-|)<\infty,
\]
where $b$ is an old bond and $c$ a final bond. 


Let \(f_{i,z}\) satisfy \cref{em:lem-affine-pointing}, and subtract
its own normalized classical pointing to obtain
\(r_{i,z}=f_{i,z}-b_i(f_i)\mathsf e_{i,z}\). Assume the classical
pointing also satisfies the hypotheses of
\cref{em:lem-affine-pointing}, has coefficient one and its total
\(S_i\) satisfies the
actual identity \(S_i(B_\xi(w))=S_k(w)\). The input local transport
estimate is the following per-anchor estimate on the composed
function \(r_{i,z}\circ B_\xi\). It has a normally convergent local
representation rooted at the cell of \(\xi z\) of norm at most
\(C\varepsilon_i r^2\xi^{4+\beta}\) for \(\beta>0\) on fixed
output radius and carrier rate. Here \(r>0\) is the fixed input
radius parameter. Each term depends on the final bond data in a
fixed-width determining collar of its connected final-cell carrier.
The norm sums the weighted suprema of these terms with the same
positive tree weight as the profile norm. Assume the same
per-anchor estimate for the companion, and for the finite-period
difference with the extra weight \(e^{\gamma P/M}\), where \(P\)
is the final shortest period and \(\gamma>0\) is fixed. The
companion and finite-period terms use the specified coherent
localizations. All constants are independent of \(i,k,P,z\).

Define the new complete pointing by
\begin{equation}\label{em:eq-anchor-transport}
 R_{i,k,y}(w)=\sum_z h_y(\xi z)\,r_{i,z}(B_\xi(w)),
\end{equation}
with periodized tents on compatible tori. Its rooted norm is at most
\(C'\varepsilon_i r^2 L^{-\beta(k-i)}\), including the stated
finite-period compatibility bound. It has zero value and first
derivative if \(r_{i,z}(0)=0\), and zero companion curvature
coefficient. Its complete pointings preserve the stated covariance.
The exact classical decomposition is
\[
 \sum_z f_{i,z}(B_\xi(w))= b_i(f_i)S_k(w)+\sum_y R_{i,k,y}(w)
\]
on the tori. For several old levels the sum of the remainders has
rooted norm at most
\[
                 \frac{C'r^2}{L^\beta-1}\sup_{i<k}\varepsilon_i.
\]
\end{proposition}

\begin{proof}
Write \(n=\xi^{-1}\). In one dimension the tent at an integral
anchor has samples \(1-|m|/n\) for \(|m|<n\) and zero otherwise,
after translating the index by the anchor times \(n\). Their sum
is \(1+2\sum_{m=1}^{n-1}(1-m/n)=n\). Products and periodization
give the exact quadrature \(\xi^4\sum_z h_y(\xi z)=1\). For
each positive weight the old root lies a bounded final distance
from \(y\). Adding a bounded join to its carrier multiplies the
norm by a fixed constant. Sum the per-anchor estimate with the
tent weights and use this quadrature. The same argument applied
to the finite-period difference gives the compatibility bound,
since the shortest final period is common to all its terms.

The old gauge invariance and covariance of the literal background
give gauge invariance of each new pointing. Every final unit
translation is an integral old translation in old units, and the
tents respect translations and signed coordinate permutations.
Finiteness for each
fixed transport suffices for the affine identity. It bounds the
response to cut-off affine data by $C_\xi(1+|\xi b_--y|)$,
uniformly in the cutoff radius, at final anchor $y$. The old
Hessian moments therefore justify the dominated cutoff limit in
the affine second derivative. To see why its kernel computes the
Hessian on these responses, first take finitely supported final
data $v$ with support $S$. The row bound gives
\[
 |(H_\xi v)(b)|\le
 \frac{C_\xi\|v\|_\infty}{1+\operatorname{dist}(\xi b_-,S)}.
\]
Thus $H_\xi v$ tends to zero at infinity. Its finite truncations
converge in supremum norm, and continuity of the old bilinear
Hessian identifies its value with the coefficient-kernel sum.
Apply this first to finite affine cutoffs and then use the stated
moment domination. A unit lattice first-difference norm gives the
same topology, since a first difference is bounded by twice the
supremum norm.

The first derivative vanishes by the chain rule and the old
vanishing first derivative. In the affine second derivative of
the composition, the second derivative of the background is
therefore multiplied by zero. The first response contributes the
old affine curvature \(\xi^2F\). The polynomially bounded gradient
disappears by the two-leg Ward identity used in
\cref{em:lem-affine-pointing}. Indeed, cut off its site field and
use the old Hessian moments to pass to the limit; those moments
hold at every fixed polynomial order. The remaining term is the old affine
test times \(\xi^4\), and that test is zero since \(b_i(r_i)=0\).
The local Hessian and response moments justify each cutoff limit.
The finite tent sum preserves this zero at each anchor.
In the direct final datum chart write the new rooted representation
as $\sum_X F_X$. A term's Hessian vanishes unless both bond
directions lie in its determining collar. That collar contains
at most $C_M(1+d_k(X))^4$ bonds, each within
$C_M(1+d_k(X))$ of the root after the bounded join. For a fixed
output ball of radius $r_{\rm out}>0$, the bidisc Cauchy estimate
therefore bounds its double Hessian moment of total degree $a$ by
$C_a r_{\rm out}^{-2}(1+d_k(X))^{8+a}\|F_X\|_\infty$.
Summing over $X$ gives every fixed moment in
\cref{em:lem-affine-pointing}, because any fixed polynomial is
bounded by a constant times the positive exponential tree weight.
This uses the weighted sum norm of the representation.
Equation \eqref{em:eq-affine-pointing} for the output now gives
zero companion coefficient.

Sum \eqref{em:eq-anchor-transport} over \(y\) and use
\(\sum_y h_y(\xi z)=1\) and the assumed exact transport of the
classical total. This gives the decomposition. Finally
\(\sum_{n\ge1}L^{-\beta n}=(L^\beta-1)^{-1}\), which gives the
claimed bound for the old levels. Holomorphy and compatibility
are preserved by the normally convergent sums on the common
output domains.
\end{proof}

Let $T_{i,k}:\mathcal B_i\to\mathcal B_k$ be bounded linear maps,
for $0\le i\le k\le N$, with
\begin{equation}\label{em:eq-transport-data}
 T_{i,i}=I,\qquad T_{j,k}T_{i,j}=T_{i,k},\qquad
 \|T_{i,k}\|\le C_T\lambda^{k-i},\qquad
 C_T\ge1,\quad 0<\lambda<1.
\end{equation}
The constants and the spaces' analytic and compatibility bounds
are independent of $N$ and of the admitted volume.
For literal transport, \cref{em:prop-anchor-transport} supplies the
decay bound whenever its per-anchor estimate is uniform in the input
profile norm and applies on these common spaces. Its fixed radius
factor is included in $C_T$. The identities below supply composition.


The semigroup requirement can be checked directly when transport is
literal pullback followed by the tent pointing of \eqref{em:eq-tents}.
If $\xi^{-1}$ is a positive integer, then
\begin{equation}\label{em:eq-tent-refinement}
 \sum_{z\in\mathbb Z^4}h_y(\xi z)h_z(t)=h_y(\xi t).
\end{equation}
In one dimension the function on the right is piecewise linear and
its knots are integers. Its linear interpolation from integer values
is therefore exact. Taking products proves the displayed identity.
Periodizing proves the same identity on compatible tori. If the
literal background pullbacks compose up to site gauge, gauge
invariance of each complete pointing removes that gauge. Applying
\eqref{em:eq-tent-refinement} to the two successive pointing sums
then gives $T_{j,k}T_{i,j}=T_{i,k}$. Absolute convergence of those
sums is required before changing their order. This argument does
not prove the decay estimate in \eqref{em:eq-transport-data}.

\begin{lemma}[A norm adapted to transport]\label{em:lem-adapted-norm}
Suppose \eqref{em:eq-transport-data} holds. Fix $\lambda<\rho<1$.
For $\psi\in\mathcal B_j$ set
\begin{equation}\label{em:eq-adapted-norm}
 \|\psi\|_{j,*}=\max_{0\le n\le N-j}
       \rho^{-n}\|T_{j,j+n}\psi\|_{\mathcal B_{j+n}}.
\end{equation}
These are Banach norms satisfying
\[
 \|\psi\|_{\mathcal B_j}\le\|\psi\|_{j,*}
 \le C_T\|\psi\|_{\mathcal B_j},\qquad
 \|T_{j,j+1}\psi\|_{j+1,*}\le\rho\|\psi\|_{j,*}
 \quad(j<N).
\]
All constants are independent of the horizon.
\end{lemma}
\begin{proof}
The term $n=0$ gives the lower bound. The transport bound gives
$\rho^{-n}\|T_{j,j+n}\psi\|\le C_T(\lambda/\rho)^n\|\psi\|$,
which proves the upper bound and completeness. For $j<N$, composition
and the change of index $r=n+1$ give
\[
 \begin{split}
 \|T_{j,j+1}\psi\|_{j+1,*}
 &=\rho\max_{1\le r\le N-j}
             \rho^{-r}\|T_{j,j+r}\psi\|_{\mathcal B_{j+r}}\\
 &\le\rho\|\psi\|_{j,*}.
 \end{split}
\]
\end{proof}

For the next result the coupling sequence is prescribed independently
of the profile recursion. Fix $m>1$, and write
$\varepsilon_{\mathrm f}(t)=t^{-1}(\log t)^m$ for $t>1$.
Let $\Theta_0,\ldots,\Theta_N>1$ be a real frame with
\begin{equation}\label{em:eq-profile-frame}
 |\Theta_j-\{\Theta_*+b_\infty(N-j)\}|\le B_a,
 \qquad b_\infty>0,\quad B_a\ge0.
\end{equation}
Here $\Theta_*$ is a lower-scale parameter that will be taken large.
The constants $b_\infty,B_a$ are fixed. For example, if
$\Theta_j=\Theta_*+\sum_{\ell=j}^{N-1}\beta_\ell$ and
$|\beta_\ell-b_\infty|\le B_b\lambda_b^\ell$ with
$0<\lambda_b<1$, then \eqref{em:eq-profile-frame} holds with
$B_a=B_b/(1-\lambda_b)$. This follows by summing the geometric tail.
No positivity of every $\beta_\ell$ is needed.

In the regular step, let $\mathsf e_j$ be the actual classical
action pointing, normalized by $b(\mathsf e_j^\infty)=1$ on the
companion. For a complete vacuum-subtracted pointing $H$, its
projection is $H-b(H^\infty)\mathsf e_j$. On the fixed-radius
spaces both this projection and the coefficient readout are assumed
bounded with a uniform constant, as supplied by the physical
coefficient estimate. They preserve the real subspaces: the readout
is real on real families and the actual classical pointing is real
on real fields. The seed $n_j$ below is the projection of
the Gaussian pointing. The map $\mathcal F_j^\perp$ is the
projection of the complete non-Gaussian formation, with the scalar
vacuum normalization retained separately. Exact transport of the
old complete pointing and this projection identify the profile
recursion below with the profile part of that step. This
identification, including the literal domains and pointings, is
part of its one-step admission.

This profile presentation must place the scalar Gaussian weight
powers entirely in the vacuum normalization. Constancy of their
total alone does not identify each complete pointing with a
constant. If the chosen pointing retains a nonconstant family
with zero total, that family must be included in
$\mathcal F_j^\perp$ and satisfy its stated bounds. Neither a
zero total nor a zero companion coefficient implies those bounds.

Fix $0<\sigma_{\mathrm p}<\sigma_{\mathrm p}^+$. At step $j<N$ suppose a jointly
holomorphic map $\mathcal F_j^{\perp}(x,\psi)\in\mathcal B_{j+1}$ is given on
$|x-\Theta_j|<\sigma_{\mathrm p}^+\Theta_j$, $\|\psi\|_{j,*}<H_*+1$.
Suppose $n_j\in\mathcal B_{j+1}$ and, on this domain,
\begin{equation}\label{em:eq-profile-formation}
 \begin{split}
 \|n_j\|_{j+1,*}&\le B_n,\\
 \|\mathcal F_j^{\perp}(x,\psi)\|_{j+1,*}
  &\le C_{H_*}\varepsilon_{\mathrm f}(\Theta_j).
 \end{split}
\end{equation}
These bounds are assumed for the same actual formation map on its
buffered domains. The buffer gives the derivative estimates used
below, by the following Cauchy argument. Frozen cutoffs, Gaussian
scaling and scalar vacuum normalization are part of that map.
A bound on an unbuffered domain alone would not give uniform
derivative estimates up to its boundary.

\begin{lemma}[Derivatives on the buffered profile domain]
\label{em:lem-profile-cauchy}
Let $X,Y$ be complex Banach spaces, $H_*\ge0$, $\Theta>0$ and
$0<\sigma_{\mathrm p}<\sigma_{\mathrm p}^+$. Suppose
$F(x,\psi)\in Y$ is holomorphic on
\[
 |x-\Theta|<\sigma_{\mathrm p}^+\Theta,\qquad
 \|\psi\|_X<H_*+1,
\]
with $\|F(x,\psi)\|_Y\le A$ there. On the inner product of balls
$|x-\Theta|\le\sigma_{\mathrm p}\Theta$, $\|\psi\|_X\le H_*$,
one has
\[
 \|D_\psi F(x,\psi)\|_{X\to Y}\le 2A,\qquad
 \|\partial_x F(x,\psi)\|_Y
   \le \frac{2A}{(\sigma_{\mathrm p}^+-\sigma_{\mathrm p})\Theta}.
\]
\end{lemma}
\begin{proof}
For any unit vector $v\in X$, the circle $|t|=1/2$ in the
variation $\psi+t v$ lies in the outer profile ball. Cauchy's
formula bounds the derivative at $t=0$ in $Y$ by $2A$.
The circle of radius $(\sigma_{\mathrm p}^+-\sigma_{\mathrm p})\Theta/2$
in the variation $x+t$ lies in the outer weight disc. The same
formula gives the second estimate. These arguments use one
complex line at a time and are independent of the dimension of
either space.
\end{proof}

\begin{theorem}[Uniform forward profile bounds]\label{em:thm-profile}
Assume \eqref{em:eq-transport-data}, \eqref{em:eq-profile-frame}
and \eqref{em:eq-profile-formation}. Choose
$H_*\ge1+2B_n/(1-\rho)$, where $\lambda<\rho<1$.
There is a floor for $\Theta_*$, depending only on the stated constants,
such that every sequence with $|x_j-\Theta_j|\le\sigma_{\mathrm p} \Theta_j$ gives a
well-defined recursion
\begin{equation}\label{em:eq-profile-recursion}
 \psi_0=0,\qquad
 \psi_{j+1}=T_{j,j+1}\psi_j+n_j+\mathcal F_j^{\perp}(x_j,\psi_j),\qquad j<N,
\end{equation}
with $\|\psi_j\|_{j,*}\le H_*$ at every $0\le j\le N$.
The solution is holomorphic on a neighborhood of this closed
coupling polydisc. For two such sequences, let
$D=\max_{0\le i<N}|x_i-\widetilde x_i|$, with $D=0$ if $N=0$.
Then
\begin{equation}\label{em:eq-profile-stability}
 \|\psi_j(x)-\psi_j(\widetilde x)\|_{j,*}
 \le C\frac{(\log \Theta_j)^m}{\Theta_j^2}D,
 \qquad 0\le j\le N.
\end{equation}
The constant and the floor are independent of the horizon and volume.
\end{theorem}
\begin{proof}
By \eqref{em:eq-profile-frame}, $\Theta_j\ge \Theta_*-B_a$.
The function $\varepsilon_{\mathrm f}(t)$ tends to zero as $t$ tends
to infinity. Choose $\Theta_*$ so that the last term in
\[
 \|\psi_{j+1}\|_{j+1,*}
 \le\rho H_*+B_n+C_{H_*}\varepsilon_{\mathrm f}(\Theta_j)
\]
is at most $(1-\rho)H_*/2$. The choice of $H_*$ makes the
right side at most $H_*$. Starting with $\psi_0=0$ proves existence
by finite induction, with every intermediate profile inside the
open formation ball. The same estimate holds on a larger open
coupling polydisc whose relative radius is strictly between
$\sigma_{\mathrm p}$ and $\sigma_{\mathrm p}^+$. Finite composition there proves holomorphy.

For stability the coupling disc and profile ball are convex.
Apply \cref{em:lem-profile-cauchy} with
$A=C_{H_*}\varepsilon_{\mathrm f}(\Theta_j)$, and integrate the
two derivatives along their straight segments. Absorb the fixed
Cauchy constants into $C_{H_*}$. With
$d_j=\|\psi_j(x)-\psi_j(\widetilde x)\|_{j,*}$ and
$k_j=(\log \Theta_j)^m/\Theta_j^2$, this yields
\[
 d_{j+1}\le
  \{\rho+C_{H_*}\varepsilon_{\mathrm f}(\Theta_j)\}d_j
  +C_{H_*}k_jD.
\]
Increase the floor until the coefficient of $d_j$ is at most
a fixed $\rho_1<1$. For $i<j$, the frame bound gives
$\Theta_i\ge \Theta_j-2B_a$. On $t>e^{m/2}$ the function
$(\log t)^m/t^2$ is decreasing. If $\Theta_i\ge \Theta_j$ this gives
$k_i\le k_j$. Otherwise choose $\Theta_*$ so $\Theta_j-2B_a\ge \Theta_j/2>1$;
then $k_i/k_j\le(\Theta_j/\Theta_i)^2\le4$.
Iterating from $d_0=0$ therefore gives
\[
 d_j\le C_{H_*}D\sum_{i<j}\rho_1^{j-1-i}k_i
 \le\frac{4C_{H_*}}{1-\rho_1}k_jD.
\]
For $N=0$ the only profile is zero, and the statement holds directly.
\end{proof}

\begin{corollary}[The scalar term of the forward profile step]
\label{em:cor-profile-scalar}
Use the admitted actual step and prescribed sequences of
\cref{em:thm-profile}. Let $\mathcal A_j$ be an ambient complex Banach space
of complete pointings on the same fixed radii, with zero value and
first derivative at the trivial field. Assume its companion readout
$H\mapsto b(H^\infty)$ and the projection
$\mathsf P_jH=H-b(H^\infty)\mathsf e_j$ from $\mathcal A_j$ to
$(\mathcal B_j,\|\cdot\|_{j,*})$ are bounded uniformly. The complete
classical pointing $\mathsf e_j\in\mathcal A_j$ has coefficient one and total
$S_j$. Let $\mathsf N_{j+1}^{\mathrm p}$ be the Gaussian pointing fixed by the reference frame of
the actual profile step in $\mathcal A_{j+1}$, with scalar vacuum factors kept separately,
and put $b^{\mathrm G}_{j+1}=b((\mathsf N_{j+1}^{\mathrm p})^\infty)$. Then
\[
 \mathsf N_{j+1}^{\mathrm p}=b^{\mathrm G}_{j+1}\mathsf e_{j+1}+n_j,
 \qquad n_j=\mathsf P_{j+1}\mathsf N_{j+1}^{\mathrm p}.
\]

Suppose the supplied unprojected non-Gaussian formation
$\mathcal F_j(x,\psi)\in\mathcal A_{j+1}$ is holomorphic on the
buffered product used in \eqref{em:eq-profile-formation} and has
norm at most $C_{H_*}\varepsilon_{\mathrm f}(\Theta_j)$ there.
It uses the same complete pointing and vacuum convention as the
actual step, with $\mathcal F_j^\perp=\mathsf P_{j+1}\mathcal F_j$.
More precisely, after the scalar vacuum factor is separated, the
actual exponent on a compatible finite torus is
\[
 -x S_{j+1}+\sum_y\bigl\{
       (T_{j,j+1}\psi)_y+(\mathsf N_{j+1}^{\mathrm p})_y
                            +(\mathcal F_j(x,\psi))_y\bigr\}.
\]
This identity is required of the supplied unprojected formation;
an identity of its projection alone would not determine the scalar.
Set
\[
    r_j(x,\psi)=b\bigl(\mathcal F_j(x,\psi)^\infty\bigr).
\]
On the inner product of balls $|x-\Theta_j|\le\sigma_{\mathrm p}\Theta_j$,
$\|\psi\|_{j,*}\le H_*$, one has
\begin{equation}\label{em:eq-profile-scalar-bounds}
  |r_j|+\|D_\psi r_j\|\le C\varepsilon_{\mathrm f}(\Theta_j),
  \qquad
  |\partial_x r_j|\le C\varepsilon_{\mathrm f}(\Theta_j)/\Theta_j.
\end{equation}
The same bounds, in the output profile norm, hold for
$\mathcal F_j^\perp$ and its derivatives. Along two of the prescribed
sequences in \cref{em:thm-profile},
\begin{equation}\label{em:eq-profile-scalar-stability}
 \big|r_j(x_j,\psi_j(x))-r_j(\widetilde x_j,\psi_j(\widetilde x))\big|
 \le C\frac{\varepsilon_{\mathrm f}(\Theta_j)}{\Theta_j}
                     \max_{0\le i<N}|x_i-\widetilde x_i|,
 \qquad j<N.
\end{equation}
The scalar action weight carried out of that step is exactly
$x_j-b^{\mathrm G}_{j+1}-r_j(x_j,\psi_j)$.
\end{corollary}

\begin{proof}
The Gaussian identity is the definition of its projection. The
bounded readout and projection give the respective outer norm
bounds from that of $\mathcal F_j$. Apply
\cref{em:lem-profile-cauchy} to each map to obtain
\eqref{em:eq-profile-scalar-bounds} and its profile counterpart.
Integrate the two derivatives along the inner weight and profile
segments. By \eqref{em:eq-profile-stability}, the profile difference
is at most $C\varepsilon_{\mathrm f}(\Theta_j)D/\Theta_j$, where $D$
is the maximum in \eqref{em:eq-profile-scalar-stability}. The
additional factor $\varepsilon_{\mathrm f}(\Theta_j)$ from $D_\psi r_j$
is bounded above on the admitted half-line. This proves
\eqref{em:eq-profile-scalar-stability}.

The old transported profile has zero companion coefficient. In
the actual exponent the classical term is $-x_j S_{j+1}$.
The Gaussian and non-Gaussian decompositions add respectively
$b^{\mathrm G}_{j+1} S_{j+1}$ and $r_j S_{j+1}$ to it. The remaining complete
pointing is exactly \eqref{em:eq-profile-recursion}. This computes
the scalar carried by one step on the prescribed input. The next
prescribed weight is independent of this computation.
\end{proof}

The transport spaces in this section have fixed physical radii.
Their norm equivalence does not remove the running-radius factor
in quadratic extraction. The theorem constructs profiles for
already prescribed couplings. Solving a backward equation for
those couplings, or admitting inherited remainder and boundary
histories, requires additional results.

\begin{thebibliography}{9}
\raggedright
\bibitem{BA} \emph{Block averaging and gauge geometry for four-dimensional lattice Yang--Mills theory}, companion manuscript, 2026.
\bibitem{GI} \emph{Gaussian block integration and free action-density correlations}, companion manuscript, 2026.
\bibitem{Bal87} T. Ba\l aban, \emph{Renormalization group approach to lattice gauge field theories. I. Generation of effective actions in a small field approximation and a coupling constant renormalization in four dimensions}, Comm. Math. Phys. \textbf{109} (1987), 249--301. \href{https://doi.org/10.1007/BF01215223}{doi:10.1007/BF01215223}.
\bibitem{Bal88c} T. Ba\l aban, \emph{Renormalization group approach to lattice gauge field theories. II. Cluster expansions}, Comm. Math. Phys. \textbf{116} (1988), 1--22. \href{https://doi.org/10.1007/BF01239022}{doi:10.1007/BF01239022}.
\end{thebibliography}
\end{document}
